% Lesson 17 — Grammar: Determiners and quantifiers: some (of), any (of), much, all (of), both (of), no, none (of), few, a few (of), little, a little (of), each, every, most (of), many, how much/how many — Anglais 1ère A — Cours signé OnBuch+
% Programme officiel MINESEC. Compiler avec : tectonic ENA18-grammar-determiners-and-quantifiers-some-of-any-of-much.tex
\newcommand{\DOCMATIERE}{Anglais}
\newcommand{\DOCNIVEAU}{1ère A}
\newcommand{\DOCMODULE}{Chapter 2 : Economic Life and Occupations}
\newcommand{\DOCLECON}{ENA18}
\newcommand{\DOCTITRE}{Lesson 17 — Grammar: Determiners and quantifiers: some (of), any (of), much, all (of), both (of), no, none (of), few, a few (of), little, a little (of), each, every, most (of), many, how much/how many}
% OnBuch+ — préambule « cours signés » ENRICHI (compilé avec tectonic / XeLaTeX)
% Système visuel « Pop » d'OnBuch+ : fond crème (jamais blanc), bordures encre
% épaisses, ombres « dures » décalées, titres Archivo Black en capitales.
% Version MATHÉMATIQUES : graphes (pgfplots/tikz), boîtes pédagogiques
% (définition, propriété, méthode, attention, le savais-tu, exercice résolu,
% exercice, TP, bilan), page de garde et signature OnBuch+ ancrée en bas.
%
% Macros attendues dans main.tex AVANT \input{preamble} :
%   \DOCMATIERE  \DOCNIVEAU  \DOCMODULE  \DOCLECON (numéro)  \DOCTITRE
\documentclass[11pt]{article}

\usepackage{fontspec}
\usepackage[english,french]{babel} % français = langue principale ; l'anglais via \eng{..} / textebox
\usepackage[a4paper,margin=2cm,top=3.4cm,bottom=2.8cm,headheight=1.4cm,footskip=1.3cm]{geometry}
\usepackage{amsmath,amssymb,amsfonts}
\usepackage{mathtools} % \xmapsto, \coloneqq... souvent utilisés par les modèles
\usepackage{cancel} % simplifications barrées (bilans de forces)
% \abs{z} (module, valeur absolue) et \norm{u} (norme), utilisés par les modèles
\providecommand{\abs}{}\let\abs\relax\DeclarePairedDelimiter{\abs}{\lvert}{\rvert}
\providecommand{\norm}{}\let\norm\relax\DeclarePairedDelimiter{\norm}{\lVert}{\rVert}
\usepackage[table]{xcolor} % [table] : fournit aussi \rowcolors (alternance de lignes)
\usepackage{pagecolor}
\usepackage{tikz}
\usetikzlibrary{arrows.meta,positioning,calc,shapes.geometric,shapes.misc,shapes.arrows,decorations.pathmorphing,decorations.pathreplacing,decorations.markings,patterns,shadows,mindmap,trees,fit,backgrounds,matrix,intersections,angles,quotes}
\usepackage{pgfplots}
\usepackage[version=4]{mhchem} % équations nucléaires \ce{^{235}_{92}U}
\usepackage[european,siunitx]{circuitikz} % schémas électriques
\pgfplotsset{compat=1.18}
\usepgfplotslibrary{fillbetween,groupplots} % aires entre courbes (calcul intégral)
\usepackage{enumitem}
\usepackage{fancyhdr}
% [most] charge toutes les bibliothèques tcolorbox (listings, minted, raster,
% documentation...) et gonfle inutilement la mémoire TeX sur un document long
% (~300 boîtes) : on ne charge que ce qui est réellement utilisé.
\usepackage[skins,breakable]{tcolorbox}
\usepackage{graphicx}
\usepackage{colortbl}
\usepackage{booktabs}
\usepackage{tabularx}
\usepackage{diagbox} % case d'en-tête diagonale des tableaux à double entrée
% Colonnes extensibles centrée / gauche / droite (C, L, R), souvent utilisées par les modèles
\newcolumntype{C}{>{\centering\arraybackslash}X}
\newcolumntype{L}{>{\raggedright\arraybackslash}X}
\newcolumntype{R}{>{\raggedleft\arraybackslash}X}
\usepackage{array}
\usepackage{multicol}
\usepackage{multirow}
\usepackage{siunitx}
% parse-numbers=false : accepte aussi \SI{2,5 \times 10^{-3}}{...}
\sisetup{locale=FR,output-decimal-marker={,},group-minimum-digits=4,parse-numbers=false}
\DeclareSIUnit\ohm{\ensuremath{\symup{\Omega}}} % U+2126 absent de la police
\DeclareSIUnit\division{div} % réglages d'oscilloscope (ms/div, V/div)
\DeclareSIUnit\tr{tr} % tours (vitesse de rotation, ex. \SI{1500}{\tr\per\minute})
\DeclareSIUnit\molar{M} % concentration molaire (ex. \SI{3}{\molar}, \SI{1}{\micro\molar})
\DeclareSIUnit\an{an} % année (le modèle confond parfois avec \year, un compteur TeX) — ex. \SI{5}{\kilo\meter\per\an}
\DeclareSIUnit\FCFA{FCFA} % franc CFA (ex. \SI{500}{\FCFA\per\kilo\gram})
\DeclareSIUnit\degreeBrix{\textdegree Brix} % teneur en sucre d'un sirop/jus (réfractométrie)
\DeclareSIUnit\month{mois} % le modèle utilise parfois \month, un compteur TeX, comme unité de durée
\usepackage{qrcode}
% \degree : commande fréquemment utilisée par les modèles pour un angle ou une
% température en dehors de \SI{}{} (non fournie par les paquets chargés).
\providecommand{\degree}{^{\circ}}
% \qed (fin de démonstration), utilisé par les modèles sans amsthm
\providecommand{\qed}{\hfill$\square$}
% \barre : double barre des valeurs interdites dans un tableau de variations (tkz-tab)
\providecommand{\barre}{\ensuremath{\|}}
% environnement proof (amsthm non chargé) utilisé par les modèles
\ifdefined\proof\else\newenvironment{proof}[1][Démonstration]{\par\noindent\textit{#1.}\ }{\qed\par}\fi
\usepackage{lastpage}
\usepackage{hyperref}
\hypersetup{hidelinks}

% ============================================================
% Palette « Pop » OnBuch+ — mêmes hex que la classe P (plus_ui.dart)
% ============================================================
\definecolor{poporange}{HTML}{F59321}
\definecolor{popdark}{HTML}{E07A0C}
\definecolor{popcream}{HTML}{FFF6E8}
\definecolor{popink}{HTML}{1C1714}
\definecolor{poppurple}{HTML}{7A5AE0}
\definecolor{popgreen}{HTML}{2E9E6A}
\definecolor{poppink}{HTML}{E0457B}
\definecolor{popblue}{HTML}{2F7DE1}
\definecolor{popgold}{HTML}{C98A12}
\definecolor{popmuted}{HTML}{8A7F76}
\definecolor{popline}{HTML}{EADFD2}
\definecolor{popgray}{HTML}{8A7F76}
\colorlet{popgrey}{popgray}
% Alias tolérants (le modèle nomme parfois les couleurs différemment)
\colorlet{popwhite}{white}
\colorlet{popblack}{popink}
\colorlet{popred}{poppink}
\colorlet{popyellow}{popgold}
% Teintes claires (fonds de tableaux/encadrés) — le modèle nomme parfois
% les couleurs avec un suffixe L (ex. popblueL) sans les avoir définies.
\colorlet{poporangeL}{poporange!20}
\colorlet{popdarkL}{popdark!20}
\colorlet{poppurpleL}{poppurple!20}
\colorlet{popgreenL}{popgreen!20}
\colorlet{poppinkL}{poppink!20}
\colorlet{popblueL}{popblue!20}
\colorlet{popgoldL}{popgold!20}
\colorlet{popredL}{popred!20}
\colorlet{popgrayL}{popgray!20}
% Teintes claires pour fonds de tableaux / zones
\colorlet{popblueL}{popblue!10!white}
\colorlet{poporangeL}{poporange!14!white}
\colorlet{popgreenL}{popgreen!12!white}
\colorlet{poppurpleL}{poppurple!12!white}
\colorlet{poppinkL}{poppink!12!white}
\colorlet{popgoldL}{popgold!14!white}

\pagecolor{popcream}

% ============================================================
% Typographie
% ============================================================
\defaultfontfeatures{Ligatures=TeX}
\setmainfont{PlusJakartaSans-Medium.ttf}[Path=../../../fonts/,BoldFont=PlusJakartaSans-Bold.ttf]
% Maths en sans-serif (Fira Math) avec chiffres et lettres droites de Plus
% Jakarta, pour que les formules se fondent dans le texte.
\usepackage[math-style=ISO,bold-style=ISO]{unicode-math}
\setmathfont{FiraMath-Regular.otf}[Path=../../../fonts/]
\setmathfont{PlusJakartaSans-Medium.ttf}[Path=../../../fonts/,range={up/num,up/{Latin,latin},\mathup}]
\usepackage{icomma} % virgule décimale sans espace en mode math
% Caractères mathématiques Unicode absents des polices de texte (Plus Jakarta,
% Archivo Black) : rendus par leur équivalent mathématique, en texte comme en
% formule — sinon ils s'affichent en carré vide (ex. « Arithmétique dans ℤ »).
\usepackage{newunicodechar}
\newunicodechar{ℝ}{\ensuremath{\mathbb{R}}}
\newunicodechar{ℤ}{\ensuremath{\mathbb{Z}}}
\newunicodechar{ℂ}{\ensuremath{\mathbb{C}}}
\newunicodechar{ℕ}{\ensuremath{\mathbb{N}}}
\newunicodechar{ℚ}{\ensuremath{\mathbb{Q}}}
\newunicodechar{𝔻}{\ensuremath{\mathbb{D}}}
\newunicodechar{α}{\ensuremath{\alpha}}
\newunicodechar{β}{\ensuremath{\beta}}
\newunicodechar{γ}{\ensuremath{\gamma}}
\newunicodechar{δ}{\ensuremath{\delta}}
\newunicodechar{ε}{\ensuremath{\varepsilon}}
\newunicodechar{θ}{\ensuremath{\theta}}
\newunicodechar{λ}{\ensuremath{\lambda}}
\newunicodechar{μ}{\ensuremath{\mu}}
\newunicodechar{σ}{\ensuremath{\sigma}}
\newunicodechar{τ}{\ensuremath{\tau}}
\newunicodechar{φ}{\ensuremath{\varphi}}
\newunicodechar{ψ}{\ensuremath{\psi}}
\newunicodechar{ω}{\ensuremath{\omega}}
\newunicodechar{Σ}{\ensuremath{\Sigma}}
% majuscules produites par \MakeUppercase dans les titres (α-aminés -> Α)
\newunicodechar{Α}{\ensuremath{\alpha}}
\newunicodechar{Β}{\ensuremath{\beta}}
\newunicodechar{Γ}{\ensuremath{\Gamma}}
\newunicodechar{Δ}{\ensuremath{\Delta}}
\newunicodechar{Ε}{\ensuremath{\varepsilon}}
\newunicodechar{Θ}{\ensuremath{\Theta}}
\newunicodechar{Λ}{\ensuremath{\Lambda}}
\newunicodechar{Μ}{\ensuremath{\mu}}
\newunicodechar{Τ}{\ensuremath{\tau}}
\newunicodechar{Φ}{\ensuremath{\Phi}}
\newunicodechar{Ψ}{\ensuremath{\Psi}}
\newunicodechar{Ω}{\ensuremath{\Omega}}
\newunicodechar{↦}{\ensuremath{\mapsto}}
\newunicodechar{∈}{\ensuremath{\in}}
\newunicodechar{∉}{\ensuremath{\notin}}
\newunicodechar{∧}{\ensuremath{\wedge}}
\newunicodechar{⇔}{\ensuremath{\Leftrightarrow}}
\newunicodechar{⟺}{\ensuremath{\Longleftrightarrow}}
\newunicodechar{⇒}{\ensuremath{\Rightarrow}}
\newunicodechar{⟹}{\ensuremath{\Longrightarrow}}
\newunicodechar{∘}{\ensuremath{\circ}}
\newunicodechar{∪}{\ensuremath{\cup}}
\newunicodechar{∩}{\ensuremath{\cap}}
\newunicodechar{≡}{\ensuremath{\equiv}}
\newunicodechar{⊂}{\ensuremath{\subset}}
\newunicodechar{⊥}{\ensuremath{\perp}}
\newunicodechar{∃}{\ensuremath{\exists}}
\newunicodechar{∀}{\ensuremath{\forall}}
\newunicodechar{∛}{\ensuremath{\sqrt[3]{\,}}}
\newunicodechar{∜}{\ensuremath{\sqrt[4]{\,}}}
\newunicodechar{⃗}{\ensuremath{{}^{\to}}}
\newunicodechar{ⁿ}{\ifmmode^{n}\else\textsuperscript{\ensuremath{n}}\fi}
\newunicodechar{ˣ}{\ifmmode^{x}\else\textsuperscript{\ensuremath{x}}\fi}
\newunicodechar{ᵇ}{\ifmmode^{b}\else\textsuperscript{\ensuremath{b}}\fi}
\newunicodechar{ʳ}{\ifmmode^{r}\else\textsuperscript{\ensuremath{r}}\fi}
\newunicodechar{ᵘ}{\ifmmode^{u}\else\textsuperscript{\ensuremath{u}}\fi}
\newunicodechar{ʸ}{\ifmmode^{y}\else\textsuperscript{\ensuremath{y}}\fi}
\newunicodechar{ᵃ}{\ifmmode^{a}\else\textsuperscript{\ensuremath{a}}\fi}
\newunicodechar{ᵉ}{\ifmmode^{e}\else\textsuperscript{\ensuremath{e}}\fi}
\newunicodechar{ᵏ}{\ifmmode^{k}\else\textsuperscript{\ensuremath{k}}\fi}
\newunicodechar{ᵖ}{\ifmmode^{p}\else\textsuperscript{\ensuremath{p}}\fi}
\newunicodechar{ᴺ}{\ifmmode^{N}\else\textsuperscript{\ensuremath{N}}\fi}
\newunicodechar{ᶜ}{\ifmmode^{c}\else\textsuperscript{\ensuremath{c}}\fi}
\newunicodechar{ʰ}{\ifmmode^{h}\else\textsuperscript{\ensuremath{h}}\fi}
\newunicodechar{ᵠ}{\ifmmode^{q}\else\textsuperscript{\ensuremath{q}}\fi}
\newunicodechar{ₙ}{\ifmmode_{n}\else\textsubscript{\ensuremath{n}}\fi}
\newunicodechar{ₐ}{\ifmmode_{a}\else\textsubscript{\ensuremath{a}}\fi}
\newunicodechar{ᵢ}{\ifmmode_{i}\else\textsubscript{\ensuremath{i}}\fi}
\newunicodechar{ᵣ}{\ifmmode_{r}\else\textsubscript{\ensuremath{r}}\fi}
\newunicodechar{ₖ}{\ifmmode_{k}\else\textsubscript{\ensuremath{k}}\fi}
\newunicodechar{ᵦ}{\ifmmode_{\beta}\else\textsubscript{\ensuremath{\beta}}\fi}
\newunicodechar{ₓ}{\ifmmode_{x}\else\textsubscript{\ensuremath{x}}\fi}
\newfontfamily\popdisplay{ArchivoBlack-Regular.ttf}[Path=../../../fonts/]
\newfontfamily\poplogo{SpaceGrotesk-Bold.ttf}[Path=../../../fonts/]
\newfontfamily\popsemi{PlusJakartaSans-SemiBold.ttf}[Path=../../../fonts/]
\newfontfamily\popxbold{PlusJakartaSans-ExtraBold.ttf}[Path=../../../fonts/]
\newfontfamily\popxboldls{PlusJakartaSans-ExtraBold.ttf}[Path=../../../fonts/,LetterSpace=3.0]

\setlist[enumerate]{leftmargin=1.7em,itemsep=5pt,topsep=4pt}
\setlist[itemize]{leftmargin=1.7em,itemsep=4pt,topsep=4pt,label={\color{poporange}\textbullet}}
\setlength{\parindent}{0pt}
\setlength{\parskip}{5pt}
\renewcommand{\arraystretch}{1.25}

% pgfplots : style Pop par défaut
\pgfplotsset{
  every axis/.append style={axis line style={popink,line width=1pt},tick style={popink},
    grid style={popline,line width=0.5pt},label style={font=\small},tick label style={font=\footnotesize}},
  popcurve/.style={poporange,line width=1.8pt,no markers,smooth},
}
% Même style disponible dans un \draw TikZ simple (hors pgfplots)
\tikzset{popcurve/.style={poporange,line width=1.8pt}}
\tikzset{popgrid/.style={popline,very thin}}
\colorlet{popcurve}{poporange} % le nom du style est parfois utilisé comme couleur

% ============================================================
% Pill / squiggle
% ============================================================
\newcommand{\poppill}[3]{%
  \tikz[baseline=-0.55ex]\node[draw=popink,line width=1.2pt,rounded corners=2.6pt,%
    fill=#2,inner xsep=6.5pt,inner ysep=3pt]{%
    \popxboldls\fontsize{7}{7}\selectfont\color{#3}\MakeUppercase{#1}};%
}
\newcommand{\popsquiggle}[1]{%
  \tikz[baseline=-0.4ex]{
    \draw[#1,line width=2.6pt,line cap=round]
      plot [smooth] coordinates {(0,0.09) (0.34,0.01) (0.68,0.21) (1.02,0.01) (1.36,0.10)};
  }%
}
% Mot-clé surligné (marqueur orange sous le texte)
\newcommand{\cle}[1]{{\popxbold\color{popdark}#1}}

% ============================================================
% En-tête / pied
% ============================================================
\pagestyle{fancy}
\fancyhf{}
\renewcommand{\headrulewidth}{0pt}
\renewcommand{\footrulewidth}{0pt}
\lhead{\raisebox{-0.15ex}{\poplogo\fontsize{13}{13}\selectfont\color{popink}OnBuch\color{poporange}+}}
\rhead{\poppill{\DOCMATIERE\ \textbullet\ \DOCNIVEAU\ \textbullet\ Leçon \DOCLECON}{white}{popink}}
\lfoot{\popxbold\fontsize{7.6}{7.6}\selectfont\color{popmuted}\MakeUppercase{Cours signé OnBuch+}\ \textbullet\ {\popsemi\DOCTITRE}}
\rfoot{\tikz[baseline=-0.6ex]\node[rounded corners=8.5pt,fill=popink,minimum height=17pt,minimum width=17pt,inner xsep=5pt,inner sep=0]{\popxbold\fontsize{8}{8}\selectfont\color{popcream}\thepage\,/\,\pageref*{LastPage}};}

% ============================================================
% Titres
% ============================================================
\newcommand{\chaptitre}[2]{%
  \begin{tcolorbox}[enhanced,colback=poporange,colframe=popink,boxrule=2.6pt,arc=11pt,
      drop shadow=popink,left=15pt,right=15pt,top=12pt,bottom=11pt,before skip=3pt,after skip=18pt]
    \poppill{#2}{popink}{popcream}\par\vspace{9pt}
    {\raggedright\hyphenpenalty=10000\exhyphenpenalty=10000\popdisplay\fontsize{22}{24}\selectfont\color{popcream}\MakeUppercase{#1}\par}
  \end{tcolorbox}
}
% Section numérotée : pastille orange avec le numéro + titre Archivo
\newcounter{popsec}
\newcommand{\coursec}[1]{%
  \stepcounter{popsec}\setcounter{popsub}{0}%
  \par\vspace{16pt}\noindent
  \begin{minipage}{\linewidth}
  \tikz[baseline=-0.7ex]\node[draw=popink,line width=1.6pt,fill=poporange,rounded corners=4pt,minimum size=22pt,inner sep=0]{\popdisplay\fontsize{12}{12}\selectfont\color{popcream}\thepopsec};%
  \hspace{8pt}{\popdisplay\fontsize{15}{17}\selectfont\color{popink}\MakeUppercase{#1}}\par
  \vspace{4pt}\noindent{\color{poporange}\rule{36pt}{3.6pt}}
  \end{minipage}\par\nopagebreak\vspace{8pt}
}
\newcounter{popsub}
\newcommand{\courssub}[1]{%
  \stepcounter{popsub}%
  \par\vspace{9pt}\noindent{\popxbold\fontsize{11.5}{13}\selectfont\color{popink}{\color{poporange}\thepopsec.\thepopsub}\hspace{6pt}#1}\par\nopagebreak\vspace{4pt}
}

% ============================================================
% Boîtes pédagogiques — même recette : fond blanc, bordure épaisse colorée,
% ombre dure encre, badge plein. Titre optionnel [#1].
% ============================================================
\tcbset{popbase/.style={breakable,enhanced,colback=white,boxrule=2.3pt,arc=9pt,
  shadow={3pt}{-3pt}{0pt}{popink},left=14pt,right=14pt,top=11pt,bottom=11pt,before skip=11pt,after skip=15pt,
  fontupper=\popsemi\fontsize{10.6}{14.5}\selectfont\color{popink},
  fontlower=\popsemi\fontsize{10.6}{14.5}\selectfont\color{popink}}}
% Coche dessinée (le glyphe ✓ n'existe pas dans les polices du document)
\newcommand{\popcheck}[1]{\tikz[baseline=-0.5ex]\draw[#1,line width=1.6pt,line cap=round,line join=round](-0.13,0.0)--(-0.04,-0.09)--(0.13,0.1);}
\providecommand{\coche}{\popcheck{popgreen}}
\providecommand{\resultat}[1]{\par\smallskip\noindent\fcolorbox{popgreen}{popgreenL}{\parbox{\dimexpr\linewidth-2\fboxsep-2\fboxrule\relax}{\textbf{Résultat.} #1}}\par\smallskip}
\newcommand{\popboxhead}[3]{\poppill{#1}{#2}{white}\ifx\relax#3\relax\else\hspace{6pt}{\popxbold\fontsize{10.6}{12}\selectfont\color{popink}#3}\fi\par\vspace{7pt}}

\newtcolorbox{aretenir}[1][]{popbase,colframe=popgreen,before upper={\popboxhead{À retenir}{popgreen}{#1}}}
% Texte en anglais (document de lecture, dialogue, extrait) : cadre
% d'étude. Un titre optionnel (source, auteur) s'affiche dans l'en-tête.
\newtcolorbox{textebox}[1][]{popbase,colframe=popblue,before upper={\popboxhead{Text}{popblue}{#1}\begin{otherlanguage}{english}},after upper={\end{otherlanguage}}}
% Marques ✓ / ✗ (forme correcte / fautive) souvent utilisées par les modèles
\providecommand{\cross}{\textcolor{poppink}{\ensuremath{\times}}}
\providecommand{\xmark}{\cross}\providecommand{\crossmark}{\cross}\providecommand{\cmark}{\ensuremath{\checkmark}}
% Ligne de réponse / justification à compléter (inventée par les modèles)
\providecommand{\justification}[1]{\par\noindent\textit{Justification~:} \dotfill\par}
% \textipa (package tipa non chargé) : la police ne couvre pas l'API -> simple passage
\providecommand{\textipa}[1]{#1}
% Soulignement (ulem non chargé) : \uline, \ul
\providecommand{\uline}[1]{\underline{#1}}\providecommand{\ul}[1]{\underline{#1}}
% \glossaire{\item ...} inventée par les modèles : liste de glossaire
\providecommand{\glossaire}[1]{\begin{itemize}#1\end{itemize}}
% \alert (beamer) inventée par les modèles : texte mis en évidence
\providecommand{\alert}[1]{\textbf{\textcolor{poppink}{#1}}}
% variantes ASCII d'ordinaux écrites en macro
\providecommand{\iere}{\textsuperscript{re}}\providecommand{\ieme}{\textsuperscript{e}}\providecommand{\ere}{\textsuperscript{re}}\providecommand{\eme}{\textsuperscript{e}}\providecommand{\er}{\textsuperscript{er}}\providecommand{\ie}{\textsuperscript{e}}
% Ordinaux français écrits en macro par les modèles : 1\ière, 3\ième
\newcommand{\ière}{\textsuperscript{re}}\newcommand{\ième}{\textsuperscript{e}}
% \exemple (inventée par les modèles) : étiquette « Exemple : »
\providecommand{\exemple}{\textbf{Exemple~:}\ }
% \square utilisable aussi hors mode math (cases à cocher)
\AtBeginDocument{\let\squareMath\square\renewcommand{\square}{\ensuremath{\squareMath}}}
% Mot ou phrase en anglais à l'intérieur d'un paragraphe français
\newcommand{\eng}[1]{\ifmmode\text{\textit{#1}}\else\begingroup\selectlanguage{english}\itshape #1\endgroup\fi}
\let\esp\eng % alias toléré
% flèches d'intonation inventées par les modèles
\providecommand{\textsearrow}{}\renewcommand{\textsearrow}{\ensuremath{\searrow}}\providecommand{\textnearrow}{}\renewcommand{\textnearrow}{\ensuremath{\nearrow}}
% \fr{..} : mot français (inventé par les modèles)
\providecommand{\fr}[1]{\textit{#1}}
\newtcolorbox{exemplebox}[1][]{popbase,colframe=poporange,before upper={\popboxhead{Exemple}{poporange}{#1}}}
\newtcolorbox{definition}[1][]{popbase,colframe=popblue,before upper={\popboxhead{Définition}{popblue}{#1}}}
\newtcolorbox{propriete}[1][]{popbase,colframe=poppurple,before upper={\popboxhead{Propriété}{poppurple}{#1}}}
\newtcolorbox{methode}[1][]{popbase,colframe=popgold,before upper={\popboxhead{Méthode}{popgold}{#1}}}
\newtcolorbox{attention}[1][]{popbase,colframe=poppink,before upper={\popboxhead{Attention}{poppink}{#1}}}
\newtcolorbox{experience}[1][]{popbase,colframe=popblue,colback=popblueL,before upper={\popboxhead{Expérience}{popblue}{#1}}}
% « Le savais-tu ? » : carte violette pleine (contexte, culture, Cameroun)
\newtcolorbox{savaistu}[1][]{popbase,colback=poppurple,colframe=popink,
  fontupper=\popsemi\fontsize{10.4}{14.2}\selectfont\color{white},
  before upper={\poppill{Le savais-tu\,?}{white}{poppurple}\ifx\relax#1\relax\else\hspace{6pt}{\popxbold\color{white}#1}\fi\par\vspace{7pt}}}
% Exercice résolu : énoncé en haut, \tcblower puis la solution sur fond crème
\newcounter{popexo}\newcounter{popexr}
\newtcolorbox{exoresolu}[1][]{popbase,colframe=poporange,colbacklower=popcream,
  segmentation style={popink,line width=1.2pt,dashed},
  before upper={\stepcounter{popexr}\popboxhead{Exercice résolu \thepopexr}{poporange}{#1}},
  before lower={\poppill{Solution}{popink}{popcream}\par\vspace{6pt}}}
% Exercice (non résolu en place) — la correction est dans la partie Corrigés
\newtcolorbox{exercice}[1][]{popbase,colframe=popink,
  before upper={\stepcounter{popexo}\popboxhead{Exercice \thepopexo}{popink}{#1}}}
% Corrigé
\newtcolorbox{corrige}[1][]{popbase,colframe=popgreen,colback=popgreenL,
  before upper={\popboxhead{Corrigé}{popgreen}{#1}}}
% Objectifs / prérequis / bilan
\newtcolorbox{objectifs}{popbase,colframe=popink,colback=poporangeL,
  before upper={\popboxhead{Objectifs de la leçon}{popink}{}}}
\newtcolorbox{prerequis}{popbase,colframe=popink,colback=popblueL,
  before upper={\popboxhead{Prérequis}{popblue}{}}}
\newtcolorbox{bilan}{popbase,colframe=popink,colback=white,boxrule=2.8pt,
  before upper={\popboxhead{Fiche bilan}{poporange}{L'essentiel à connaître pour l'examen}}}
% Figure encadrée avec légende
\newtcolorbox{popfigure}[1][]{enhanced,breakable=false,colback=white,colframe=popline,boxrule=1.4pt,arc=8pt,
  left=10pt,right=10pt,top=10pt,bottom=8pt,before skip=10pt,after skip=12pt,halign=center,
  lower separated=false,halign lower=center,
  fontlower=\popsemi\fontsize{9}{11.5}\selectfont\color{popmuted},
  #1}
\newcommand{\legende}[1]{\par\vspace{6pt}{\popsemi\fontsize{9}{11.5}\selectfont\color{popmuted}#1}}

% ============================================================
% Signature OnBuch+
% ============================================================
\newcommand{\poplogoinline}[1]{{\poplogo\fontsize{#1}{#1}\selectfont\color{popink}OnBuch\color{poporange}+}}
\newcommand{\signatureonbuch}{%
  \begin{tcolorbox}[enhanced,colback=popink,colframe=popink,boxrule=2.2pt,arc=10pt,
      drop shadow=poporange,left=14pt,right=14pt,top=10pt,bottom=10pt,before skip=0pt,after skip=0pt]
    \tikz[baseline=-0.7ex]\node[circle,fill=popgreen,draw=popcream,line width=1.3pt,minimum size=22pt,inner sep=0]{\popcheck{white}};%
    \hspace{9pt}\begin{minipage}[c]{0.62\linewidth}
      {\popxbold\fontsize{12}{13}\selectfont\color{popcream}Signé\ \textbullet\ {\poplogo OnBuch\color{poporange}+}}\par\vspace{2pt}
      {\raggedright\popsemi\fontsize{8}{10.5}\selectfont\color{popline}\MakeUppercase{Équipe pédagogique OnBuch+}\\\MakeUppercase{Conforme au programme officiel MINESEC}\par}
    \end{minipage}\hfill
    {\poplogo\fontsize{22}{22}\selectfont\color{popcream}OnBuch\color{poporange}+}
  \end{tcolorbox}%
}

% ============================================================
% Page de garde
% #1 titre · #2 matière · #3 niveau · #4 module · #5 n° leçon · #6 durée ·
% #7 description / objectifs (texte ou itemize)
% ============================================================
\newcommand{\pagegarde}[7]{%
  \thispagestyle{empty}
  \newgeometry{a4paper,margin=1.7cm,top=1.6cm,bottom=1.4cm}%
  \begin{tikzpicture}[remember picture,overlay]
    \fill[poporange!18] ([xshift=-3.2cm,yshift=-3.6cm]current page.north east) circle (4.6cm);
    \fill[poppurple!14] ([xshift=2.4cm,yshift=7.6cm]current page.south west) circle (3.8cm);
  \end{tikzpicture}%
  {\poplogo\fontsize{30}{30}\selectfont\color{popink}OnBuch\color{poporange}+}\hfill
  \raisebox{0.6ex}{\poppill{Cours signé \textbullet\ Premium}{popink}{popcream}}\par
  \vspace{1pt}\popsquiggle{poppurple}\par
  \vspace{8pt}
  \begin{tcolorbox}[enhanced,colback=poporange,colframe=popink,boxrule=2.8pt,arc=13pt,
      drop shadow=popink,left=18pt,right=18pt,top=12pt,bottom=13pt,after skip=10pt]
    \poppill{#2 \textbullet\ #3}{popink}{popcream}\hspace{5pt}\poppill{#4}{white}{popink}\par\vspace{8pt}
    {\popdisplay\fontsize{14}{14}\selectfont\color{popink}LEÇON #5}\par\vspace{4pt}
    {\raggedright\hyphenpenalty=10000\exhyphenpenalty=10000\popdisplay\fontsize{25}{27}\selectfont\color{popcream}\MakeUppercase{#1}\par}
    \vspace{8pt}
    {\popxbold\fontsize{9.5}{9.5}\selectfont\color{popink}\MakeUppercase{Durée conseillée : #6}}
  \end{tcolorbox}
  \begin{tcolorbox}[enhanced,colback=white,colframe=popink,boxrule=2.2pt,arc=10pt,
      drop shadow=popink,left=16pt,right=16pt,top=10pt,bottom=10pt,after skip=8pt]
    \poppill{Dans cette leçon}{poppurple}{white}\par\vspace{5pt}
    \popsemi\fontsize{9.3}{12.6}\selectfont\color{popink}#7
  \end{tcolorbox}
  \noindent\begin{minipage}[t]{0.487\linewidth}
    \begin{tcolorbox}[enhanced,colback=popgreen,colframe=popink,boxrule=2.2pt,arc=10pt,
        drop shadow=popink,left=11pt,right=11pt,top=10pt,bottom=10pt,height=3.1cm,valign=center]
      \tcbox[colback=white,colframe=popink,boxrule=1.3pt,arc=5pt,boxsep=0pt,nobeforeafter,
        left=3pt,right=3pt,top=3pt,bottom=3pt,valign=center]{\qrcode[height=1.9cm]{https://play.google.com/store/apps/details?id=com.luvvix.onbuch&hl=fr}}%
      \hspace{8pt}\begin{minipage}[c]{0.5\linewidth}
        {\popdisplay\fontsize{11}{12.5}\selectfont\color{white}\MakeUppercase{Télécharge OnBuch}}\par\vspace{4pt}
        {\raggedright\popsemi\fontsize{8.4}{11}\selectfont\color{white}Gratuit \textbullet\ Google Play\\Tuteur IA, annales, concours, résultats\par}
      \end{minipage}
    \end{tcolorbox}
  \end{minipage}\hfill
  \begin{minipage}[t]{0.487\linewidth}
    \begin{tcolorbox}[enhanced,colback=poppurple,colframe=popink,boxrule=2.2pt,arc=10pt,
        drop shadow=popink,left=11pt,right=11pt,top=10pt,bottom=10pt,height=3.1cm,valign=center]
      \tikz[baseline=-0.7ex]\node[circle,fill=white,draw=popink,line width=1.3pt,minimum size=1.6cm,inner sep=0]{\popdisplay\fontsize{24}{24}\selectfont\color{poppurple}+};%
      \hspace{8pt}\begin{minipage}[c]{0.6\linewidth}
        {\popdisplay\fontsize{11}{12.5}\selectfont\color{white}\MakeUppercase{Passe à OnBuch+}}\par\vspace{4pt}
        {\raggedright\popsemi\fontsize{8.4}{11}\selectfont\color{white}Tous les cours signés, Tuteur IA illimité, fiches PDF — onglet OnBuch+ de l'app\par}
      \end{minipage}
    \end{tcolorbox}
  \end{minipage}
  \vspace{8pt}
  \popcontenu
  \vfill
  \signatureonbuch
  \restoregeometry
  \newpage
}

% Rangée « Ce cours contient » (page de garde)
\newcommand{\poptile}[3]{% #1 couleur #2 gros texte #3 légende
  \begin{tcolorbox}[enhanced,colback=#1,colframe=popink,boxrule=2pt,arc=9pt,shadow={2.5pt}{-2.5pt}{0pt}{popink},
      left=6pt,right=6pt,top=9pt,bottom=9pt,halign=center,valign=center,height=1.75cm,width=\linewidth,nobeforeafter]
    {\popdisplay\fontsize{12.5}{13}\selectfont\color{white}\MakeUppercase{#2}}\par\vspace{3pt}
    {\centering\popsemi\fontsize{7.8}{9.5}\selectfont\color{white}#3\par}
  \end{tcolorbox}}
\newcommand{\popcontenu}{%
  \par\noindent{\popxboldls\fontsize{7.5}{7.5}\selectfont\color{popmuted}CE COURS SIGNÉ CONTIENT}\par\vspace{7pt}
  \noindent\begin{minipage}[t]{0.235\linewidth}\poptile{popblue}{Cours}{complet, illustré, pas à pas}\end{minipage}\hfill
  \begin{minipage}[t]{0.235\linewidth}\poptile{poporange}{Méthodes}{et exercices résolus}\end{minipage}\hfill
  \begin{minipage}[t]{0.235\linewidth}\poptile{poppink}{Exercices}{type examen + corrigés}\end{minipage}\hfill
  \begin{minipage}[t]{0.235\linewidth}\poptile{popgreen}{Fiche bilan}{l'essentiel à retenir}\end{minipage}\par}

% Sommaire
\newtcolorbox{sommaire}{enhanced,colback=white,colframe=popink,boxrule=2.2pt,arc=10pt,
  drop shadow=popink,left=18pt,right=18pt,top=13pt,bottom=13pt,before skip=6pt,after skip=20pt,
  before upper={\poppill{Sommaire}{poppurple}{white}\par\vspace{9pt}\popsemi\fontsize{10.6}{14.5}\selectfont\color{popink}}}

% Signature de fin de cours — poussée tout en bas de la dernière page
\newcommand{\findecours}{%
  \par\vfill\nopagebreak
  \begin{center}\popsquiggle{poporange}\end{center}\vspace{4pt}
  \signatureonbuch
}
\AtBeginDocument{\def\llbracket{\mathopen{[\mkern-3.5mu[}}\def\rrbracket{\mathclose{]\mkern-3.5mu]}}} % intervalles d'entiers (glyphe ⟦ absent de la police)
% Glyphes absents de Fira Math : redéfinis à partir de glyphes disponibles
\AtBeginDocument{%
  \def\setminus{\mathbin{\backslash}}\let\smallsetminus\setminus
  \def\vdots{\mathord{\vbox{\baselineskip3.2pt\lineskiplimit0pt\kern4pt\hbox{.}\hbox{.}\hbox{.}}}}%
  \def\ddots{\mathinner{\mkern1mu\raise6.5pt\hbox{.}\mkern2mu\raise3.6pt\hbox{.}\mkern2mu\raise0.7pt\hbox{.}\mkern1mu}}%
  \def\triangle{\mathord{\scalebox{0.9}{$\symup{\Delta}$}}}%
  \def\nabla{\mathord{\raisebox{\depth}{\scalebox{1}[-1]{$\symup{\Delta}$}}}}%
  \def\checkmark{\popcheck{popdark}}%
  \def\star{\mathbin{\ast}}\def\bigstar{\mathord{\ast}}%
  \def\diamond{\mathbin{\scalebox{0.6}{$\Diamond$}}}%
  \def\bigcup{\mathop{\mathchoice{\raisebox{-0.3em}{\scalebox{1.7}{$\cup$}}}{\scalebox{1.25}{$\cup$}}{\cup}{\cup}}\displaylimits}%
  \def\bigcap{\mathop{\mathchoice{\raisebox{-0.3em}{\scalebox{1.7}{$\cap$}}}{\scalebox{1.25}{$\cap$}}{\cap}{\cap}}\displaylimits}%
  \def\lesseqqgtr{\mathrel{\lessgtr}}\def\bigoplus{\mathop{\oplus}\displaylimits}%
}
\providecommand{\ssi}{si et seulement si} % abréviation fréquente
\AtBeginDocument{\providecommand{\wideparen}{\overparen}} % arc de cercle

%% FIGFIX-BEGIN v4 — correctifs globaux des figures (docs/onbuch-generation/tools/figfix)
\usepackage{adjustbox}\usepackage{etoolbox}
% toute figure TikZ/pgfplots plus large que la ligne est réduite à la largeur disponible
\BeforeBeginEnvironment{tikzpicture}{\begin{adjustbox}{max width=\linewidth}}
\AfterEndEnvironment{tikzpicture}{\end{adjustbox}}
% légendes pgfplots lisibles par-dessus les courbes
\pgfplotsset{every axis/.append style={legend style={fill=white,fill opacity=0.92,text opacity=1,draw=black!15,inner sep=3pt}}}
% étiquettes d'arêtes (syntaxe edge["..."]) avec fond blanc
\tikzset{every edge quotes/.append style={fill=white,inner sep=1.2pt}}
%% FIGFIX-END
\begin{document}

\pagegarde{Lesson 17 — Grammar: Determiners and quantifiers: some (of), any (of), much, all (of), both (of), no, none (of), few, a few (of), little, a little (of), each, every, most (of), many, how much/how many}{Anglais}{1ère A}{Chapter 2 : Economic Life and Occupations}{ENA18}{2 h}{Cette leçon de grammaire présente les déterminants et quantifieurs anglais (some, any, much, many, all, both, no, none, few, a few, little, a little, each, every, most) ainsi que les questions how much / how many. À travers l'observation d'exemples contextualisés, des règles illustrées et des exercices progressifs, les élèves apprennent à exprimer la quantité sans les erreurs typiques des francophones.
\vspace{4pt}
\begin{multicols}{2}\small
\begin{itemize}
\item À la fin de cette leçon, je sais distinguer les noms comptables des noms incomptables pour choisir le bon quantifieur
\item À la fin de cette leçon, je sais employer some et any selon le type de phrase (affirmative, négative, interrogative, offre)
\item À la fin de cette leçon, je sais utiliser much, many, a lot of et poser des questions avec how much / how many
\item À la fin de cette leçon, je sais exprimer la totalité avec all, both, each, every et most
\item À la fin de cette leçon, je sais exprimer la petite quantité et l'absence avec few, a few, little, a little, no et none
\item À la fin de cette leçon, je sais employer la structure « quantifieur + of + the/my/these + nom » correctement
\end{itemize}\end{multicols}}

\begin{sommaire}
\begin{multicols}{2}
\begin{enumerate}
\item Observation et découverte guidée
\item Some, any, no, none
\item Much, many, how much / how many
\item Few, a few, little, a little
\item All, both, each, every, most
\item Synthèse, comparaison avec le français et entraînement
\item Activité d'intégration
\item Méthodes et exercices résolus
\item Exercices
\item Corrigés des exercices
\item Fiche bilan
\end{enumerate}
\end{multicols}
\end{sommaire}

\begin{objectifs}
\begin{itemize}
\item À la fin de cette leçon, je sais distinguer les noms comptables des noms incomptables pour choisir le bon quantifieur
\item À la fin de cette leçon, je sais employer some et any selon le type de phrase (affirmative, négative, interrogative, offre)
\item À la fin de cette leçon, je sais utiliser much, many, a lot of et poser des questions avec how much / how many
\item À la fin de cette leçon, je sais exprimer la totalité avec all, both, each, every et most
\item À la fin de cette leçon, je sais exprimer la petite quantité et l'absence avec few, a few, little, a little, no et none
\item À la fin de cette leçon, je sais employer la structure « quantifieur + of + the/my/these + nom » correctement
\item À la fin de cette leçon, je sais éviter les erreurs fréquentes des francophones (much people, any dans les affirmatives, no + not...)
\item À la fin de cette leçon, je sais utiliser ces quantifieurs dans des phrases et un court texte sur la vie économique
\end{itemize}
\end{objectifs}

\begin{prerequis}
\begin{itemize}
\item Distinction nom comptable / nom incomptable (Seconde)
\item Articles a/an, the et article zéro (Seconde)
\item Formation des phrases négatives et interrogatives avec les auxiliaires do/does/did
\item Pluriel des noms et noms invariables (money, information, news)
\item Notion de déterminant vue en début d'année (possessifs, démonstratifs)
\end{itemize}
\end{prerequis}

\newpage

\coursec{Observation et découverte guidée}

Au marché Mokolo de Yaoundé, Amina veut acheter des provisions pour la semaine. Elle discute avec son amie en anglais :

\eng{“How much rice do you need? — Not much, just a few kilos. And how many tomatoes? — A few, but no onions, there are none left!”}

En quelques secondes, Amina a utilisé six quantifieurs différents : \cle{how much}, \cle{not much}, \cle{a few}, \cle{how many}, \cle{no}, \cle{none}. Et voici le problème : en anglais, on ne peut pas dire \eng{I have much friends} ni \eng{I don't have some money}. Chaque déterminant a sa place, selon que le nom est \textbf{comptable} ou \textbf{incomptable}, et selon que la phrase est \textbf{affirmative}, \textbf{négative} ou \textbf{interrogative}. Cette leçon vous donne les clés pour quantifier comme un anglophone.

\textbf{Problématique :} comment choisir le bon quantifieur devant un nom anglais ?

\courssub{Un dialogue au marché : première observation}

\begin{textebox}[At the market — dialogue d'observation]
— How much oil do we need?

— Not much, just a little. And how many eggs?

— A few. There are no tomatoes left, but there are some onions. Most traders accept Mobile Money; both shops are open, and every customer gets a gift.

— Perfect! And how much sugar do you want?

— Just a little, please. There isn't much juice left at home, but there's some water.
\end{textebox}

\begin{definition}[Glossaire du dialogue]
\begin{itemize}
\item \cle{oil} (n. incomptable) : huile
\item \cle{eggs} (n. comptable pluriel) : œufs
\item \cle{left} (participe) : restant, qui reste
\item \cle{trader} (n.) : commerçant, vendeur
\item \cle{customer} (n.) : client
\item \cle{juice} (n. incomptable) : jus
\end{itemize}
\end{definition}

Lisez ce dialogue à voix haute avec votre voisin, en respectant l'intonation montante des questions en \eng{how much / how many}. Puis faites l'activité de repérage suivante.

\begin{experience}[Repérage guidé]
\begin{enumerate}
\item \textbf{Soulignez} tous les mots placés juste devant les noms du dialogue : \eng{much oil}, \eng{a little}, \eng{many eggs}, \eng{a few}, \eng{no tomatoes}, \eng{some onions}, \eng{most traders}, \eng{both shops}, \eng{every customer}, \eng{much sugar}, \eng{much juice}, \eng{some water}.
\item \textbf{Classez les noms} en deux colonnes :
\begin{center}
\begin{tabularx}{\linewidth}{|X|X|}
\hline
\rowcolor{popblueL} \textbf{Countable nouns (comptables)} & \textbf{Uncountable nouns (incomptables)} \\
\hline
eggs, tomatoes, onions, traders, shops, customers, gifts & oil, sugar, juice, water, money \\
\hline
\end{tabularx}
\end{center}
\item \textbf{Observez} : quels quantifieurs apparaissent dans une \textbf{question} ? dans une \textbf{négation} ? dans une \textbf{affirmation} ?
\end{enumerate}
\end{experience}

\begin{corrige}[Repérage guidé]
\begin{itemize}
\item Dans les \textbf{questions} : \eng{how much oil}, \eng{how many eggs}, \eng{how much sugar}.
\item Dans les \textbf{négations} : \eng{not much}, \eng{no tomatoes}, \eng{isn't much juice}.
\item Dans les \textbf{affirmations} : \eng{a little}, \eng{a few}, \eng{some onions}, \eng{most traders}, \eng{both shops}, \eng{every customer}, \eng{some water}.
\end{itemize}
Première constatation : \cle{much} apparaît surtout dans les questions et les négations, et devant des noms \textbf{incomptables} ; \cle{many} devant des noms \textbf{comptables} ; \cle{some} dans les affirmations.
\end{corrige}

\courssub{Qu'est-ce qu'un déterminant quantifieur ?}

\begin{definition}[Quantifier]
A \cle{quantifier} is a word used before a noun (or alone, as a pronoun) to express a quantity — \emph{how much} or \emph{how many}. Un \textbf{quantifieur} est un mot placé devant un nom (ou employé seul, comme pronom) pour indiquer une quantité précise ou vague.
\end{definition}

En français, on dit « du riz », « des œufs », « peu d'espoir », « quelques amis » : l'article partitif ou indéfini fait une partie du travail. En anglais, il n'existe pas d'équivalent direct de « du / de la / des » : c'est le quantifieur qui porte toute l'information de quantité. D'où l'importance capitale de bien le choisir.

\courssub{La grande question : comptable ou incomptable ?}

Avant tout quantifieur, l'anglophone se pose une question que le francophone oublie souvent : \textbf{ce nom peut-il être compté ?}

\begin{propriete}[Noms comptables et incomptables]
\begin{itemize}
\item Un nom \textbf{comptable} (\cle{countable noun}) a un singulier et un pluriel : \eng{an egg, two eggs} ; \eng{a shop, both shops}. On peut le précéder de \eng{a/an} et d'un nombre.
\item Un nom \textbf{incomptable} (\cle{uncountable noun}) n'a pas de pluriel et ne prend jamais \eng{a/an} : \eng{oil, water, money, rice, information, advice, furniture}. Pour le compter, on utilise une unité : \eng{a bottle of oil}, \eng{a kilo of rice}, \eng{a piece of advice}.
\end{itemize}
\end{propriete}

\begin{attention}[Incomptables piégeux pour les francophones]
Certains noms incomptables en anglais sont comptables en français : \eng{information} (une information), \eng{advice} (un conseil), \eng{news} (une nouvelle), \eng{money} (de l'argent), \eng{bread} (un pain), \eng{furniture} (un meuble), \eng{homework} (un devoir). On dit donc \eng{much information}, jamais \eng{many informations}.
\end{attention}

\begin{popfigure}
\begin{tikzpicture}[mindmap, grow cyclic, every node/.style={concept, align=center, font=\small},
root concept/.append style={concept color=popblue, font=\bfseries},
level 1/.append style={level distance=3.6cm, sibling angle=120},
level 2/.append style={level distance=2.4cm, sibling angle=50}]
\node[root concept]{Quantifiers}
  child[concept color=popgreen] { node[align=center]{countable\\ nouns}
    child { node[concept color=popgreenL, text=popdark, align=center]{many\\ few\\ a few} }
    child { node[concept color=popgreenL, text=popdark, align=center]{how many\\ each\\ every} }
  }
  child[concept color=poporange] { node[align=center]{uncountable\\ nouns}
    child { node[concept color=poporangeL, text=popdark, align=center]{much\\ little\\ a little} }
    child { node[concept color=poporangeL, text=popdark] {how much} }
  }
  child[concept color=poppurple] { node[align=center]{both\\ (les deux)}
    child { node[concept color=poppurpleL, text=popdark, align=center]{some, any\\ no, none} }
    child { node[concept color=poppurpleL, text=popdark, align=center]{all, most\\ both} }
  };
\end{tikzpicture}
\legende{Carte mentale : le choix du quantifieur dépend d'abord de la nature du nom.}
\end{popfigure}

\courssub{Question-guide : où apparaît chaque quantifieur ?}

Reprenons nos observations du dialogue et formulons des \textbf{hypothèses} avant la règle officielle.

\begin{experience}[Hypothèses des élèves]
En groupes, complétez ces hypothèses à partir du dialogue :
\begin{enumerate}
\item \eng{Some} semble s'employer dans les phrases \ldots\ (affirmatives / négatives / interrogatives) ?
\item \eng{Much} et \eng{many} semblent s'employer surtout dans les phrases \ldots\ ?
\item \eng{No} contient-il déjà une négation ? Peut-on dire \eng{there aren't no tomatoes} ?
\item \eng{A few} s'emploie devant les noms \ldots, \eng{a little} devant les noms \ldots
\end{enumerate}
\end{experience}

\begin{corrige}[Hypothèses]
\begin{enumerate}
\item \eng{Some} apparaît dans les phrases \textbf{affirmatives} : \eng{there are some onions}, \eng{there's some water}.
\item \eng{Much} et \eng{many} apparaissent surtout dans les \textbf{questions} et les \textbf{négations} : \eng{how much oil?}, \eng{not much}, \eng{there isn't much juice}.
\item \eng{No} est déjà négatif : \eng{there are no tomatoes} = \eng{there aren't any tomatoes}. La double négation \eng{there aren't no tomatoes} est une faute grave en anglais standard.
\item \eng{A few} + nom \textbf{comptable pluriel} (\eng{a few eggs}) ; \eng{a little} + nom \textbf{incomptable} (\eng{a little oil}).
\end{enumerate}
\end{corrige}

\begin{exemplebox}[Exemples traduits]
\begin{itemize}
\item \eng{There are a few eggs.} « Il reste quelques œufs. »
\item \eng{There is little hope.} « Il y a peu d'espoir. »
\item \eng{There is a little water.} « Il y a un peu d'eau. »
\item \eng{There are no tomatoes left.} « Il ne reste plus de tomates. »
\item \eng{How many customers came today?} « Combien de clients sont venus aujourd'hui ? »
\item \eng{How much money do you have?} « Combien d'argent as-tu ? »
\end{itemize}
\end{exemplebox}

\courssub{Premier tableau de synthèse : qui va avec qui ?}

\begin{center}
\begin{tabularx}{\linewidth}{|l|X|X|X|}
\hline
\rowcolor{popblueL} \textbf{Quantifieur} & \textbf{Type de nom} & \textbf{Type de phrase} & \textbf{Français} \\
\hline
some & comptable pluriel ou incomptable & affirmative & du, des, certains \\
\hline
any & comptable pluriel ou incomptable & négative, interrogative & de, du, aucun \\
\hline
no & comptable ou incomptable & affirmative à sens négatif & aucun, pas de \\
\hline
much & incomptable & négative, interrogative & beaucoup de \\
\hline
many & comptable pluriel & négative, interrogative & beaucoup de \\
\hline
a few & comptable pluriel & affirmative & quelques \\
\hline
few & comptable pluriel & affirmative à sens négatif & peu de \\
\hline
a little & incomptable & affirmative & un peu de \\
\hline
little & incomptable & affirmative à sens négatif & peu de \\
\hline
all, most & comptable pluriel ou incomptable & affirmative & tout, la plupart \\
\hline
both & comptable pluriel (deux) & affirmative & les deux \\
\hline
each, every & comptable singulier & affirmative & chaque \\
\hline
\end{tabularx}
\end{center}

\begin{methode}[Choisir le bon quantifieur en trois questions]
\begin{enumerate}
\item \textbf{Le nom est-il comptable ?} Si oui → \eng{many, few, a few, each, every, both}. Si non → \eng{much, little, a little}.
\item \textbf{La phrase est-elle affirmative, négative ou interrogative ?} Affirmative → \eng{some} ; négative ou interrogative → \eng{any} (ou \eng{much/many}).
\item \textbf{Le sens est-il positif ou négatif ?} \eng{A few} = « quelques » (sens positif : il y en a) ; \eng{few} = « presque pas » (sens négatif : il n'y en a guère). Même contraste entre \eng{a little} et \eng{little}.
\end{enumerate}
\end{methode}

\begin{exoresolu}[Premier exercice de découverte]
Énoncé : complétez avec le quantifieur qui convient (\eng{much, many, a few, a little, some, no}).
\begin{enumerate}
\item How \ldots{} rice do you cook every day?
\item I have \ldots{} friends in Bamenda; we meet every Saturday.
\item There is \ldots{} sugar left; go and buy a kilo at Mokolo market.
\item She added \ldots{} salt to the soup — just a small quantity.
\item There are \ldots{} taxis at this hour; the station is empty.
\item Would you like \ldots{} water?
\end{enumerate}
\tcblower
Solution détaillée :
\begin{enumerate}
\item \eng{much} — \eng{rice} est incomptable et la phrase est interrogative. « Combien de riz cuisines-tu chaque jour ? »
\item \eng{a few} — \eng{friends} est comptable pluriel, phrase affirmative, sens positif (« quelques amis »).
\item \eng{no} — le sens est négatif (« plus de sucre ») ; on pourrait aussi dire \eng{isn't any sugar}.
\item \eng{a little} — \eng{salt} est incomptable, sens positif (« un peu de sel »).
\item \eng{no} — la gare est vide : il n'y a aucun taxi. (\eng{There are no taxis} = \eng{There aren't any taxis}.)
\item \eng{some} — dans une offre ou une invitation (\eng{Would you like...?}), on garde \eng{some} même en question, car on attend une réponse positive. « Veux-tu de l'eau ? »
\end{enumerate}
\end{exoresolu}

\begin{attention}[Pièges immédiats des francophones]
\begin{itemize}
\item Ne dites jamais \eng{I have much friends} : \eng{friends} est comptable → \eng{I have many friends} ou mieux, à l'affirmatif, \eng{I have a lot of friends}.
\item Ne dites jamais \eng{I don't have some money} : en négation, \eng{some} devient \eng{any} → \eng{I don't have any money}.
\item \eng{How many} exige un pluriel : \eng{how many egg} est faux → \eng{how many eggs}.
\item Pas de double négation : \eng{there isn't no oil} est incorrect → \eng{there is no oil} ou \eng{there isn't any oil}.
\end{itemize}
\end{attention}

\begin{aretenir}[Ce qu'il faut retenir de cette observation]
\begin{itemize}
\item Un quantifieur indique une quantité devant un nom (ou seul comme pronom : \eng{none left}).
\item La première question à se poser : le nom est-il \textbf{comptable} (\eng{many, a few}) ou \textbf{incomptable} (\eng{much, a little}) ?
\item La deuxième question : la phrase est-elle \textbf{affirmative} (\eng{some}), \textbf{négative} ou \textbf{interrogative} (\eng{any, much, many}) ?
\item La troisième question : le sens est-il \textbf{positif} (\eng{a few, a little}) ou \textbf{négatif} (\eng{few, little, no, none}) ?
\end{itemize}
\end{aretenir}

\begin{savaistu}[Pourquoi l'anglais distingue-t-il much et many ?]
Le vieil anglais possédait déjà deux mots distincts : \emph{micel} (grande quantité, devenu \eng{much}) et \emph{manig} (grand nombre, devenu \eng{many}). Le français, lui, utilise « beaucoup de » pour les deux situations — c'est pourquoi les francophones confondent si souvent \eng{much} et \eng{many}. Au Cameroun anglophone, au marché de Buea ou de Bamenda, vous entendrez chaque jour : \eng{How much is a bucket of garri?} et \eng{How many cups do you want?} — la distinction est vivante dans la langue de tous les jours.
\end{savaistu}

\begin{exercice}[À vous de jouer]
\begin{enumerate}
\item Classez ces noms en deux colonnes (comptable / incomptable) : \eng{money, banana, oil, chair, information, egg, water, book, rice, customer}.
\item Transformez : \eng{There are some onions.} → forme négative avec \eng{any}, puis forme négative avec \eng{no}.
\item Traduisez : « Combien de tomates as-tu achetées ? — Quelques-unes. Et combien d'huile ? — Pas beaucoup, juste un peu. »
\end{enumerate}
\end{exercice}

\begin{corrige}[À vous de jouer]
\begin{enumerate}
\item Comptables : \eng{banana, chair, egg, book, customer}. Incomptables : \eng{money, oil, information, water, rice}.
\item \eng{There aren't any onions.} / \eng{There are no onions.}
\item \eng{How many tomatoes did you buy? — A few. And how much oil? — Not much, just a little.}
\end{enumerate}
\end{corrige}

Nous avons posé les fondations : la distinction comptable / incomptable et le trio affirmation–négation–question. Dans la prochaine section, nous étudierons en détail le système \cle{some / any / no / none}, le cœur de cette leçon.

\coursec{Some, any, no, none}

Dans la section précédente, nous avons observé le dialogue au marché de Mokolo et repéré les mots \eng{some}, \eng{any}, \eng{no} et \eng{none} dans des phrases authentiques. Nous allons maintenant dégager la règle générale : ces petits mots, appelés \cle{déterminants} et \cle{quantifieurs}, obéissent en anglais à une logique très régulière — mais différente du français. C'est l'une des sources d'erreurs les plus fréquentes chez les francophones.

\courssub{Observation : retour sur le texte}

Relisons attentivement ces phrases extraites de notre dialogue d'observation :

\begin{exemplebox}[Phrases observées]
\begin{itemize}
\item \eng{I need some tomatoes and some oil.} « J'ai besoin de tomates et d'huile. » (phrase affirmative)
\item \eng{Sorry, I don't have any tomatoes today.} « Désolé, je n'ai pas de tomates aujourd'hui. » (phrase négative)
\item \eng{Do you have any change?} « Avez-vous de la monnaie ? » (question)
\item \eng{Would you like some groundnuts?} « Voulez-vous des arachides ? » (question, mais c'est une \textbf{offre})
\item \eng{There is no pepper left.} « Il n'y a plus de piment. » (négation avec \eng{no})
\item \eng{How many customers came this morning? — None.} « Combien de clients sont venus ce matin ? — Aucun. » (\eng{none} seul, sans nom)
\end{itemize}
\end{exemplebox}

\textbf{Questions de découverte :} dans quelles phrases trouve-t-on \eng{some} ? Dans quelles phrases trouve-t-on \eng{any} ? Que remarque-t-on dans la question \eng{Would you like some groundnuts?} ? La règle va confirmer vos observations.

\courssub{La règle générale : some, any, no}

\begin{propriete}[Règle fondamentale : some / any / no]
\begin{enumerate}
\item \cle{Some} s'emploie dans les phrases \textbf{affirmatives} : \eng{I have some money.} « J'ai de l'argent. »
\item \cle{Any} s'emploie dans les phrases \textbf{négatives} (avec \eng{not}) et dans les \textbf{questions} : \eng{I don't have any money.} « Je n'ai pas d'argent. » ; \eng{Do you have any money?} « As-tu de l'argent ? »
\item \textbf{Exception importante} : dans les questions qui sont des \textbf{offres} ou des \textbf{demandes polies} (on attend une réponse positive), on garde \cle{some} : \eng{Would you like some tea?} « Voudrais-tu du thé ? » ; \eng{Can I have some water, please?} « Puis-je avoir de l'eau, s'il vous plaît ? »
\item \cle{No} + nom exprime la négation \textbf{sans} \eng{not} : \eng{There is no water.} = \eng{There isn't any water.} « Il n'y a pas d'eau. » On n'utilise jamais \eng{no} et \eng{not} ensemble.
\end{enumerate}
\end{propriete}

\begin{popfigure}
\begin{center}
\begin{tikzpicture}[
  box/.style={draw=popblue, fill=popblueL, rounded corners, align=center, minimum width=3.8cm, minimum height=1.1cm, font=\small},
  lab/.style={font=\small\bfseries, text=popdark}
]
\node[lab] (aff) at (0,2.2) {Phrase affirmative};
\node[lab] (neg) at (0,0.6) {Phrase négative};
\node[lab] (int) at (0,-1.0) {Question};
\node[box, align=center] (b1) at (6,2.2){\textbf{some}\\ I have some friends.};
\node[box, align=center] (b2) at (6,0.6){\textbf{any} (avec not) ou \textbf{no}\\ I don't have any money. / I have no money.};
\node[box, fill=poporangeL, draw=poporange, align=center] (b3) at (6,-1.0){\textbf{any} en général\\ Do you have any change?};
\node[box, fill=popgreenL, draw=popgreen, align=center] (b4) at (6,-2.4){mais \textbf{some} si offre / demande polie\\ Would you like some tea?};
\draw[-{Stealth}, thick, popblue] (aff) -- (b1);
\draw[-{Stealth}, thick, popblue] (neg) -- (b2);
\draw[-{Stealth}, thick, poporange] (int) -- (b3);
\draw[-{Stealth}, thick, popgreen, dashed] (b3) -- (b4) node[midway, right, font=\scriptsize\itshape, text=popdark]{exception};
\end{tikzpicture}
\end{center}
\legende{Le choix entre some, any et no selon le type de phrase.}
\end{popfigure}

\courssub{Some et any : devant quels noms ?}

\eng{Some} et \eng{any} se placent devant un \textbf{nom pluriel} ou un \textbf{nom indénombrable} (uncountable). Ils correspondent en français à « du, de la, des » ou à l'absence d'article après une négation.

\begin{center}
\begin{tabularx}{\linewidth}{|X|X|X|}
\hline
\rowcolor{popblueL} English & Français & Type de nom \\
\hline
\eng{some water} & de l'eau & indénombrable \\
\hline
\eng{some books} & des livres & pluriel \\
\hline
\eng{some money} & de l'argent & indénombrable \\
\hline
\eng{any information} & de l'information / aucune information & indénombrable \\
\hline
\eng{any students} & des élèves / d'élèves & pluriel \\
\hline
\end{tabularx}
\end{center}

\begin{attention}[Noms indénombrables : piège majeur]
En anglais, \eng{money}, \eng{water}, \eng{bread}, \eng{information}, \eng{news}, \eng{advice}, \eng{furniture} sont \textbf{indénombrables} : ils n'ont pas de pluriel. On dit \eng{some information}, jamais \emph{some informations}. On dit \eng{some advice} (« des conseils »), jamais \emph{some advices}. Pour compter, on utilise \eng{a piece of} : \eng{a piece of advice} « un conseil », \eng{two pieces of information} « deux informations ».
\end{attention}

\begin{exemplebox}[Exemples traduits]
\begin{itemize}
\item \eng{I have some friends in Buea.} « J'ai des amis à Buéa. »
\item \eng{She bought some rice at the market.} « Elle a acheté du riz au marché. »
\item \eng{I don't have any money.} « Je n'ai pas d'argent. »
\item \eng{There aren't any taxis at this hour.} « Il n'y a pas de taxis à cette heure-ci. »
\item \eng{Do you have any brothers or sisters?} « As-tu des frères et sœurs ? »
\item \eng{Would you like some coffee?} « Voudrais-tu du café ? » (offre → \eng{some})
\item \eng{Could I have some sugar, please?} « Pourrais-je avoir du sucre, s'il vous plaît ? » (demande polie → \eng{some})
\end{itemize}
\end{exemplebox}

\begin{attention}[Offres et demandes polies : on garde some]
Les francophones appliquent souvent mécaniquement « question = any ». Faux ! Quand la question est une \textbf{offre} ou une \textbf{demande} où l'on espère une réponse positive, on utilise \eng{some} : \eng{Would you like some coffee?} « Veux-tu du café ? » ; \eng{Can I have some money?} « Puis-je avoir de l'argent ? ». Compare : \eng{Is there any coffee left?} « Reste-t-il du café ? » (simple information → \eng{any}).
\end{attention}

\courssub{Any au sens de « n'importe quel »}

Dans les phrases \textbf{affirmatives}, \eng{any} existe aussi, mais avec un sens particulier : « n'importe quel(le)(s) », « peu importe lequel ».

\begin{definition}[Any = n'importe quel]
Dans une phrase affirmative, \cle{any} signifie « n'importe quel », sans restriction de choix. Il peut se mettre devant un singulier, un pluriel ou un indénombrable.
\end{definition}

\begin{exemplebox}[Exemples traduits]
\begin{itemize}
\item \eng{You can take any bus to Douala.} « Tu peux prendre n'importe quel bus pour Douala. »
\item \eng{Any student can answer this question.} « N'importe quel élève peut répondre à cette question. »
\item \eng{Come and see me any day you like.} « Viens me voir n'importe quel jour. »
\item \eng{You can ask me any question.} « Tu peux me poser n'importe quelle question. »
\end{itemize}
\end{exemplebox}

\courssub{No : la négation du nom}

\begin{propriete}[No + nom]
\cle{No} se place directement devant un nom (singulier, pluriel ou indénombrable) et rend la phrase négative \textbf{sans} utiliser \eng{not}. Deux formulations équivalentes :
\begin{itemize}
\item \eng{There is no water.} = \eng{There isn't any water.} « Il n'y a pas d'eau. »
\item \eng{I have no money.} = \eng{I don't have any money.} « Je n'ai pas d'argent. »
\item \eng{There are no buses on Sundays.} = \eng{There aren't any buses on Sundays.} « Il n'y a pas de bus le dimanche. »
\end{itemize}
Avec \eng{no}, le verbe reste à la forme \textbf{affirmative}. La double négation est interdite en anglais standard.
\end{propriete}

\begin{attention}[Jamais de double négation]
L'erreur la plus fréquente des francophones (et des apprenants influencés par le parler populaire) : \emph{I don't have no money.} Cette phrase est \textbf{incorrecte} en anglais standard : deux négations s'annulent ou choquent. Il faut choisir : \eng{I don't have any money.} ou \eng{I have no money.} De même : \eng{I didn't see anybody.} ou \eng{I saw nobody.}, jamais \emph{I didn't see nobody.}
\end{attention}

\courssub{None : le pronom}

\begin{propriete}[None]
\cle{None} est un \textbf{pronom} : il remplace le nom et s'emploie donc \textbf{seul}, jamais devant un nom. Il répond souvent aux questions en \eng{How many...?} ou \eng{How much...?} et signifie « aucun(e) », « pas un(e) seul(e) ».
\begin{itemize}
\item \eng{How many students came? — None.} « Combien d'élèves sont venus ? — Aucun. »
\item \eng{How much petrol is left? — None.} « Combien reste-t-il d'essence ? — Plus du tout. »
\end{itemize}
Pour préciser, on utilise \cle{none of} + \eng{the / these / my...} + nom (ou pronom) : \eng{None of my friends came.} « Aucun de mes amis n'est venu. » ; \eng{None of them spoke English.} « Aucun d'entre eux ne parlait anglais. »
\end{propriete}

\begin{attention}[None ou no ?]
Devant un nom, utilisez \eng{no} ; seul ou avec \eng{of}, utilisez \eng{none}. On dit \eng{no students} mais \eng{none of the students}. Jamais \emph{none students} ni \emph{no of the students}.
\end{attention}

\courssub{Some of, any of : « certains de... »}

\begin{propriete}[Some of / any of + déterminant + nom]
Pour dire « certains de... », « quelques-uns de... », on utilise \cle{some of} (ou \cle{any of}) suivi \textbf{obligatoirement} d'un déterminant (\eng{the}, \eng{my}, \eng{these}...) puis du nom, ou d'un pronom (\eng{us}, \eng{them}).
\begin{itemize}
\item \eng{Some of my friends live in Bamenda.} « Certains de mes amis vivent à Bamenda. »
\item \eng{Some of the students passed the exam.} « Certains des élèves ont réussi l'examen. »
\item \eng{I don't know any of these traders.} « Je ne connais aucun de ces commerçants. »
\item \eng{Did any of them come?} « Est-ce que certains d'entre eux sont venus ? »
\end{itemize}
Attention : sans \eng{of}, pas de déterminant : \eng{some friends} mais \eng{some of my friends}. Jamais \emph{some of friends}.
\end{propriete}

\courssub{Les composés : somebody, anybody, nobody, something, anything, nothing}

Les composés en \eng{-body/-one} (personnes) et \eng{-thing} (choses) suivent exactement la même logique que \eng{some/any/no} :

\begin{center}
\begin{tabularx}{\linewidth}{|X|X|X|X|}
\hline
\rowcolor{popblueL} & Affirmatif & Négatif / Question & Négation (verbe affirmatif) \\
\hline
Personnes & \eng{somebody / someone} & \eng{anybody / anyone} & \eng{nobody / no one} \\
\hline
Choses & \eng{something} & \eng{anything} & \eng{nothing} \\
\hline
Lieux & \eng{somewhere} & \eng{anywhere} & \eng{nowhere} \\
\hline
\end{tabularx}
\end{center}

\begin{exemplebox}[Exemples traduits]
\begin{itemize}
\item \eng{Somebody stole my phone at the market.} « Quelqu'un a volé mon téléphone au marché. »
\item \eng{Did anyone call me?} « Quelqu'un m'a-t-il appelé ? »
\item \eng{I didn't see anybody.} = \eng{I saw nobody.} « Je n'ai vu personne. »
\item \eng{I have something to tell you.} « J'ai quelque chose à te dire. »
\item \eng{She said nothing.} = \eng{She didn't say anything.} « Elle n'a rien dit. »
\item \eng{There is nowhere to park in Douala.} « Il n'y a nulle part où se garer à Douala. »
\end{itemize}
\end{exemplebox}

\begin{savaistu}[No one s'écrit en deux mots]
Contrairement à \eng{someone} et \eng{anyone}, qui s'écrivent en un seul mot, \eng{no one} s'écrit toujours en \textbf{deux mots} : \eng{No one came to the meeting.} « Personne n'est venu à la réunion. » C'est une particularité orthographique de l'anglais, souvent sanctionnée dans les rédactions. Retenez aussi que tous ces composés sont \textbf{singuliers} : \eng{Somebody \textbf{is} waiting.}, \eng{Nobody \textbf{knows} the answer.}
\end{savaistu}

\courssub{Comparaison avec le français}

\begin{center}
\begin{tabularx}{\linewidth}{|X|X|X|}
\hline
\rowcolor{popblueL} Français & English & Remarque \\
\hline
J'ai des amis. & \eng{I have some friends.} & \eng{some} souvent omis en anglais parlé : \eng{I have friends.} \\
\hline
Je n'ai pas d'argent. & \eng{I don't have any money. / I have no money.} & deux traductions possibles \\
\hline
As-tu de la monnaie ? & \eng{Do you have any change?} & question → \eng{any} \\
\hline
Tu veux du thé ? & \eng{Would you like some tea?} & offre → \eng{some} \\
\hline
Aucun élève n'est venu. & \eng{No student came. / None of the students came.} & verbe affirmatif en anglais \\
\hline
Je n'ai rien vu. & \eng{I saw nothing. / I didn't see anything.} & jamais de double négation \\
\hline
\end{tabularx}
\end{center}

\courssub{Exercice résolu et entraînement}

\begin{exoresolu}[Complétez avec some, any, no ou none]
Énoncé — Complete the sentences with \emph{some}, \emph{any}, \emph{no} or \emph{none}.
\begin{enumerate}
\item There is \dots\ bread left; go and buy some at the bakery.
\item Would you like \dots\ palm wine?
\item I don't have \dots\ brothers, but I have two sisters.
\item How many passengers were on the bus? — \dots . It was empty.
\item You can ask \dots\ question you like during the lesson.
\item \dots\ of my classmates come from Bafoussam.
\end{enumerate}
\tcblower
Solution détaillée :
\begin{enumerate}
\item \eng{no} — phrase négative sans \eng{not} : « Il ne reste plus de pain » (= \eng{There isn't any bread}).
\item \eng{some} — question qui est une \textbf{offre} (\eng{Would you like...?}) : on garde \eng{some}.
\item \eng{any} — phrase négative avec \eng{don't} : \eng{I don't have any brothers}.
\item \eng{None} — réponse à \eng{How many...?}, pronom seul : « Aucun. »
\item \eng{any} — phrase affirmative au sens de « n'importe quelle » : \eng{any question you like}.
\item \eng{Some} — « certains de mes camarades » : \eng{Some of my classmates} (avec le déterminant \eng{my}).
\end{enumerate}
\end{exoresolu}

\begin{exercice}[À vous de jouer]
\begin{enumerate}
\item Traduisez en anglais : a) « Il n'y a pas d'eau dans le robinet. » (deux façons) ; b) « Voudriez-vous du jus ? » ; c) « Aucun de mes amis n'a de voiture. » ; d) « Tu peux prendre n'importe quel bendskin pour aller au marché. » ; e) « Je n'ai vu personne au stade. » (deux façons).
\item Corrigez les erreurs : a) \emph{I don't have no money.} b) \emph{Would you like any tea?} c) \emph{None students came.} d) \emph{Some of students passed.}
\end{enumerate}
\end{exercice}

\begin{corrige}[Exercice « À vous de jouer »]
\begin{enumerate}
\item a) \eng{There is no water in the tap. / There isn't any water in the tap.} b) \eng{Would you like some juice?} c) \eng{None of my friends have a car. / None of my friends has a car.} d) \eng{You can take any bendskin to go to the market.} e) \eng{I saw nobody at the stadium. / I didn't see anybody at the stadium.}
\item a) \eng{I don't have any money.} ou \eng{I have no money.} (double négation interdite) ; b) \eng{Would you like some tea?} (offre) ; c) \eng{No students came.} ou \eng{None of the students came.} ; d) \eng{Some of the students passed.} (il manque le déterminant après \eng{of}).
\end{enumerate}
\end{corrige}

\begin{aretenir}[L'essentiel de la section]
\begin{itemize}
\item \eng{Some} : affirmatives + offres et demandes polies (\eng{Would you like some...?}).
\item \eng{Any} : négatives (avec \eng{not}) et questions ; sens de « n'importe quel » dans les affirmatives.
\item \eng{No} + nom, verbe affirmatif ; \eng{none} seul ou \eng{none of the/my/these} + nom.
\item \eng{Some of} / \eng{any of} exigent un déterminant : \eng{some of my friends}.
\item Composés : \eng{somebody, anybody, nobody, something, anything, nothing} — même logique ; \eng{no one} en deux mots.
\item Jamais de double négation : \eng{I don't have any}, pas \emph{I don't have no}.
\end{itemize}
\end{aretenir}

\coursec{Much, many, how much / how many}

Dans la section précédente, nous avons vu que \eng{some}, \eng{any}, \eng{no} et \eng{none} servent à exprimer une quantité indéfinie. Mais comment dire « beaucoup de » ? Comment demander « combien ? ». L'anglais impose ici un choix que le français ne fait pas : il faut d'abord se demander si le nom est \cle{comptable} (on peut le compter : \eng{one student, two students}) ou \cle{incomptable} (on ne peut pas le compter : \eng{money, water, time}). De ce choix dépendent \eng{much} ou \eng{many}, et \eng{how much} ou \eng{how many}.

\courssub{Observation : un dialogue au marché de Mokolo}

\begin{textebox}[At Mokolo market, Yaoundé]
\textbf{Customer:} Good morning, madam. How much is this bag of rice?

\textbf{Trader:} This one? It is 12 000 francs. We don't have many bags left, you know. The season has been difficult.

\textbf{Customer:} Really? I don't have much money with me today. Can you reduce the price?

\textbf{Trader:} Sorry, my sister. I bought it at a high price myself. But I can give you a lot of extra tomatoes for free.

\textbf{Customer:} How many tomatoes?

\textbf{Trader:} As many as you want! I have too many tomatoes this week and not much time to sell them before they go bad.

\textbf{Customer:} All right, I will take the rice and the tomatoes. How much altogether?

\textbf{Trader:} 12 500 francs. Thank you very much!
\end{textebox}

\textbf{Glossaire :} \emph{trader} (n.) : commerçante ; \emph{left} (adj.) : restant ; \emph{to reduce} (v.) : réduire, baisser ; \emph{for free} (loc. adv.) : gratuitement ; \emph{to go bad} (v.) : se gâter, pourrir ; \emph{altogether} (adv.) : au total.

\textbf{Questions de compréhension :}
\begin{enumerate}
\item How much does the bag of rice cost?
\item Why does the trader give extra tomatoes?
\item Does the customer have much money?
\end{enumerate}

\begin{corrige}[Questions de compréhension]
\begin{enumerate}
\item \eng{It costs 12 000 francs.} « Il coûte 12 000 francs. »
\item \eng{Because she has too many tomatoes and not much time to sell them before they go bad.} « Parce qu'elle a trop de tomates et peu de temps pour les vendre avant qu'elles ne se gâtent. »
\item \eng{No, he doesn't have much money with him today.} « Non, il n'a pas beaucoup d'argent sur lui aujourd'hui. »
\end{enumerate}
\end{corrige}

\begin{aretenir}[Ce que nous observons]
Dans ce dialogue, \eng{much} apparaît devant \eng{money} et \eng{time} (noms incomptables), tandis que \eng{many} apparaît devant \eng{bags} et \eng{tomatoes} (noms pluriels comptables). Remarquons aussi : \eng{how much} pour demander le prix, \eng{how many} pour demander un nombre, et \eng{a lot of} à la forme affirmative.
\end{aretenir}

\courssub{La règle fondamentale : much ou many ?}

\begin{propriete}[Règle de much et many]
\begin{itemize}
\item \cle{much} + \textbf{nom incomptable} (singulier, sans \emph{s}) : \eng{much money, much time, much water, much information, much patience}.
\item \cle{many} + \textbf{nom pluriel comptable} : \eng{many students, many cars, many questions, many people}.
\item \eng{much} et \eng{many} s'emploient surtout dans les phrases \textbf{négatives} et \textbf{interrogatives}. À la forme \textbf{affirmative}, l'anglais préfère généralement \eng{a lot of} ou \eng{lots of}, qui conviennent aux deux types de noms.
\end{itemize}
\end{propriete}

\begin{exemplebox}[Exemples traduits]
\begin{itemize}
\item \eng{We don't have much time.} « Nous n'avons pas beaucoup de temps. »
\item \eng{Do you drink much water?} « Bois-tu beaucoup d'eau ? »
\item \eng{There aren't many pupils in the classroom today.} « Il n'y a pas beaucoup d'élèves dans la classe aujourd'hui. »
\item \eng{Did you ask many questions?} « As-tu posé beaucoup de questions ? »
\item \eng{She has a lot of friends.} « Elle a beaucoup d'amis. » (affirmatif, comptable)
\item \eng{She has a lot of patience.} « Elle a beaucoup de patience. » (affirmatif, incomptable)
\item \eng{They earn lots of money.} « Ils gagnent beaucoup d'argent. »
\end{itemize}
\end{exemplebox}

\begin{attention}[Much à l'affirmatif]
\eng{much} et \eng{many} ne sont pas interdits à l'affirmatif, mais ils y sonnent lourds ou formels : \eng{I have much work} est correct mais peu naturel ; on dira plutôt \eng{I have a lot of work}. En revanche, après \eng{too}, \eng{so}, \eng{as} et \eng{very}, \eng{much/many} sont obligatoires même à l'affirmatif : \eng{I have so much work!}, \eng{There are too many cars in Douala.}
\end{attention}

\courssub{How much ou how many ? Demander la quantité, le prix, le nombre}

\begin{propriete}[How much / how many]
\begin{itemize}
\item \cle{how much} + nom incomptable : pour demander une \textbf{quantité} ou un \textbf{prix}. \eng{How much time do we have?} ; \eng{How much is this fabric?}
\item \cle{how many} + nom pluriel comptable : pour demander un \textbf{nombre}. \eng{How many pupils are there in your class?}
\item \eng{How much} peut aussi s'employer \textbf{seul}, sans nom, pour demander le prix : \eng{How much is it?} ou simplement \eng{How much?}
\end{itemize}
\end{propriete}

\begin{exemplebox}[Exemples traduits]
\begin{itemize}
\item \eng{How much is this bag of rice?} « Combien coûte ce sac de riz ? »
\item \eng{How much money do you need?} « De combien d'argent as-tu besoin ? »
\item \eng{How many brothers do you have?} « Combien de frères as-tu ? »
\item \eng{How many bendskins are there in Bamenda?} « Combien y a-t-il de motos-taxis à Bamenda ? »
\end{itemize}
\end{exemplebox}

\begin{popfigure}
\begin{tikzpicture}[
  node distance=1.2cm and 2.2cm,
  q/.style={rectangle, rounded corners, draw=popdark, fill=popblueL, font=\bfseries, align=center, minimum height=0.9cm},
  br/.style={rectangle, rounded corners, draw=popline, fill=popcream, align=center, minimum height=1.1cm, text width=5.6cm},
  arr/.style={-{Stealth[length=3mm]}, thick, popdark}
]
\node[q] (how) {How \ldots ?};
\node[br, below left=of how, align=center] (much){\textbf{much} + uncountable\\ price or quantity\\ \emph{How much is this shirt?}\\ \emph{How much time is left?}};
\node[br, below right=of how, align=center] (many){\textbf{many} + countable plural\\ number\\ \emph{How many students?}\\ \emph{How many books?}};
\draw[arr] (how) -- (much);
\draw[arr] (how) -- (many);
\end{tikzpicture}
\legende{Choisir entre \eng{how much} et \eng{how many} : prix/quantité contre nombre.}
\end{popfigure}

\begin{methode}[Répondre à how much / how many : les réponses courtes]
En anglais parlé, on répond souvent par une réponse courte, sans répéter le nom :
\begin{enumerate}
\item \eng{How much money do you have?} — \eng{Not much.} / \eng{A lot.} / \eng{None.}
\item \eng{How many children are there?} — \eng{Not many.} / \eng{A lot.} / \eng{None.}
\item \eng{How much is it?} — \eng{Two thousand francs.}
\end{enumerate}
On peut aussi donner la réponse complète : \eng{I don't have much money.}, \eng{There aren't many children.}
\end{methode}

\courssub{Too much, too many, so much, so many}

\begin{propriete}[Intensité : too / so + much / many]
Même logique comptable/incomptable :
\begin{itemize}
\item \eng{too much} + incomptable = « trop de » : \eng{There is too much noise in the market.} « Il y a trop de bruit au marché. »
\item \eng{too many} + pluriel = « trop de » : \eng{Too many people were at the stadium.} « Trop de gens étaient au stade. »
\item \eng{so much} / \eng{so many} = « tellement de » : \eng{She has so much energy!} ; \eng{There were so many cars!}
\end{itemize}
\end{propriete}

\courssub{Tableau récapitulatif et comparaison avec le français}

\begin{center}
\begin{tabularx}{\linewidth}{|X|X|X|}
\hline
\rowcolor{popblueL} Français & English & Remarque \\
\hline
beaucoup d'argent & much money / a lot of money & incomptable \\
\hline
beaucoup d'amis & many friends / a lot of friends & pluriel comptable \\
\hline
Combien coûte ce tissu ? & How much is this fabric? & prix = how much \\
\hline
Combien de frères as-tu ? & How many brothers do you have? & nombre = how many \\
\hline
pas beaucoup de temps & not much time & incomptable \\
\hline
pas beaucoup de voitures & not many cars & pluriel \\
\hline
trop de bruit / trop de gens & too much noise / too many people & too + much/many \\
\hline
\end{tabularx}
\end{center}

En français, « beaucoup de » et « combien de » fonctionnent avec tous les noms. En anglais, il faut d'abord classer le nom : c'est le réflexe essentiel à acquérir.

\begin{exoresolu}[Exercice : much ou many ?]
Complétez avec \eng{much} ou \eng{many} : 1. How \ldots{} water do you drink every day? 2. There aren't \ldots{} taxis at night. 3. I don't have \ldots{} homework today. 4. How \ldots{} goals did Eto'o score? 5. We don't have \ldots{} information about the exam.
\tcblower
\textbf{Solution détaillée :} 1. \eng{much} — \eng{water} est incomptable. 2. \eng{many} — \eng{taxis} est un pluriel comptable. 3. \eng{much} — \eng{homework} est incomptable en anglais (jamais de \emph{s}). 4. \eng{many} — \eng{goals} est comptable. 5. \eng{much} — \eng{information} est incomptable en anglais, contrairement au français « des informations ».
\end{exoresolu}

\begin{attention}[Erreurs fréquentes des francophones]
\begin{itemize}
\item \eng{much people} est \textbf{impossible} : \eng{people} est pluriel → \eng{many people}.
\item \eng{how many money} est \textbf{impossible} : \eng{money} est incomptable → \eng{how much money}.
\item Attention aux noms français comptables mais anglais incomptables : \eng{information}, \eng{news}, \eng{advice}, \eng{homework}, \eng{money}, \eng{bread}. On dit \eng{much information}, jamais \eng{many informations}.
\item Ne traduisez pas « beaucoup » par \eng{much} à l'affirmatif : préférez \eng{a lot of} : \eng{He has a lot of money.}
\end{itemize}
\end{attention}

\begin{savaistu}[« How much? » au Royaume-Uni]
Dans un magasin britannique, il suffit de montrer un article et de demander simplement \eng{How much?} pour connaître son prix : la phrase complète \eng{How much is it?} est sous-entendue. C'est l'une des premières phrases que tout voyageur apprend, avec la réponse typique du vendeur : \eng{It's five pounds.}
\end{savaistu}

\begin{exercice}[À vous de jouer]
1. Traduisez : « Combien de temps avons-nous ? » ; « Il y a trop de voitures à Douala. » ; « Elle n'a pas beaucoup d'argent. » 2. Posez trois questions avec \eng{how much} ou \eng{how many} à un camarade sur sa famille, son école et son argent de poche, puis répondez par des réponses courtes (\eng{Not much}, \eng{A lot}, \eng{None}).
\end{exercice}

\begin{corrige}[À vous de jouer]
1. \eng{How much time do we have?} ; \eng{There are too many cars in Douala.} ; \eng{She doesn't have much money.} 2. Réponses libres ; exemples : \eng{How many brothers and sisters do you have?} ; \eng{How many pupils are there in your class?} ; \eng{How much pocket money do you get every week?}
\end{corrige}

\coursec{Few, a few, little, a little}

Dans la section précédente, nous avons étudié \cle{much} et \cle{many}, les quantificateurs de la grande quantité, ainsi que les questions \eng{how much?} et \eng{how many?}. Nous descendons maintenant à l'autre extrémité de l'échelle : comment exprimer la \textbf{petite quantité} en anglais ? C'est là que se cache l'un des pièges les plus subtils pour les francophones : l'anglais distingue \eng{few} de \eng{a few}, et \eng{little} de \eng{a little}. La présence ou l'absence du petit mot \eng{a} change complètement le sens — et l'émotion — de la phrase.

\courssub{Observation : un même marché, deux points de vue}

\begin{textebox}[At the Mokolo market — two opinions]
Two friends, Amina and Brenda, are talking about business at Mokolo market in Yaoundé.

\emph{Amina:} How was business today?

\emph{Brenda:} Terrible. Few customers came to my shop this morning, and I made little money. I am really worried.

\emph{Amina:} Don't give up! A few customers bought my tomatoes, and I still have a little time before the market closes. There is a little hope for both of us!

\emph{Brenda:} You are right. Quite a few traders told me that business improves in the afternoon.

\emph{Amina:} Exactly. And a few of my regular customers always come after four o'clock.
\end{textebox}

\textbf{Glossaire :} \emph{customer} (nom, client) ; \emph{worried} (adjectif, inquiet) ; \emph{to give up} (verbe, abandonner) ; \emph{hope} (nom, espoir) ; \emph{trader} (nom, commerçant) ; \emph{regular} (adjectif, habituel).

\textbf{Questions de compréhension :}
\begin{enumerate}
\item Why is Brenda worried?
\item Who is optimistic, Amina or Brenda? Which words show it?
\item What do traders say about the afternoon?
\end{enumerate}

Observons attentivement le dialogue. Brenda, découragée, dit : \eng{Few customers came} et \eng{I made little money}. Amina, optimiste, dit : \eng{A few customers bought my tomatoes} et \eng{I still have a little time}. Les quantités sont probablement les mêmes, mais la \textbf{vision} est différente : sans \eng{a}, le message est négatif ; avec \eng{a}, il est positif.

\courssub{La règle d'or : avec ou sans « a »}

\begin{propriete}[La nuance fondamentale : « a » change tout]
\begin{itemize}
\item \cle{few} + nom \textbf{pluriel comptable} = « peu de », sens \textbf{négatif} (presque pas, moins qu'espéré, déception).
\item \cle{a few} + nom \textbf{pluriel comptable} = « quelques », sens \textbf{positif} (un petit nombre, mais c'est déjà bien).
\item \cle{little} + nom \textbf{incomptable} = « peu de », sens \textbf{négatif} (presque pas, insuffisant).
\item \cle{a little} + nom \textbf{incomptable} = « un peu de », sens \textbf{positif} (une petite quantité, mais suffisante).
\end{itemize}
\textbf{Moyen mnémotechnique :} le \eng{a} est un « plus » : il ajoute de l'optimisme. Sans \eng{a}, on souligne ce qui \emph{manque} ; avec \eng{a}, on souligne ce qui \emph{existe}.
\end{propriete}

\begin{exemplebox}[Exemples traduits : la nuance en action]
\begin{itemize}
\item \eng{Few people came to the meeting.} « Peu de gens sont venus à la réunion. » (déception : on en attendait davantage)
\item \eng{A few people came to the meeting.} « Quelques personnes sont venues à la réunion. » (satisfaction : ce n'est pas rien)
\item \eng{There is little milk left.} « Il reste peu de lait. » (inquiétude : il va en manquer)
\item \eng{There is a little milk left.} « Il reste un peu de lait. » (soulagement : assez pour le thé)
\item \eng{Few students passed the test.} « Peu d'élèves ont réussi le test. » (c'est dommage)
\item \eng{I have a few friends in Limbe.} « J'ai quelques amis à Limbé. » (c'est déjà bien)
\item \eng{There is little hope.} « Il y a peu d'espoir. »
\item \eng{I speak a little English.} « Je parle un peu anglais. »
\end{itemize}
\end{exemplebox}

\begin{attention}[Erreur fréquente des francophones]
\eng{A little} ne s'emploie \textbf{jamais} avec un nom comptable pluriel, et \eng{a few} jamais avec un incomptable.
\begin{itemize}
\item Faux : \eng{a little friends} $\rightarrow$ Correct : \eng{a few friends} (quelques amis).
\item Faux : \eng{a few money} $\rightarrow$ Correct : \eng{a little money} (un peu d'argent).
\item Faux : \eng{few water} $\rightarrow$ Correct : \eng{little water} (peu d'eau).
\end{itemize}
Retenez : \eng{few / a few} comptent des \textbf{unités} (books, chairs, students), \eng{little / a little} mesurent une \textbf{masse} (water, time, money, hope).
\end{attention}

\courssub{Tableau récapitulatif}

\begin{center}
\begin{tabularx}{\linewidth}{|l|l|X|X|}
\hline
\rowcolor{popblueL} Quantificateur & Type de nom & Sens français & Exemple \\
\hline
few & pluriel comptable & peu de (négatif) & \eng{Few taxis stop here.} \\
\hline
a few & pluriel comptable & quelques (positif) & \eng{A few taxis stop here.} \\
\hline
little & incomptable & peu de (négatif) & \eng{We have little time.} \\
\hline
a little & incomptable & un peu de (positif) & \eng{We have a little time.} \\
\hline
quite a few & pluriel comptable & pas mal de & \eng{Quite a few traders use Mobile Money.} \\
\hline
\end{tabularx}
\end{center}

\courssub{Frise des quantités : situer few et little}

\begin{popfigure}
\begin{tikzpicture}[scale=1]
\draw[-{Stealth}, thick, popdark] (0,0) -- (12.5,0);
\foreach \x/\txt/\col in {0.4/{none}/popink, 2.6/{little / few\\(presque rien)}/poporange, 5.0/{a little / a few\\(un peu)}/popgold, 7.4/{some}/popgreen, 9.8/{a lot of /\\much, many}/popblue, 12.0/{all}/poppurple} {
  \fill[\col] (\x,0) circle (0.09);
  \node[above, align=center, font=\small, text=\col] at (\x,0.1) {\txt};
}
\node[below, font=\small\itshape, text=popmuted] at (6.2,-0.35) {quantité croissante →};
\end{tikzpicture}
\legende{De « rien du tout » à « tout » : \eng{few/little} expriment la quasi-absence, \eng{a few/a little} la petite quantité positive.}
\end{popfigure}

\courssub{Quite a few : le complément Bac}

\begin{definition}[\cle{Quite a few}]
\eng{Quite a few} + nom pluriel comptable signifie \emph{a fairly large number} : « pas mal de », « un bon nombre de ». Attention au paradoxe : bien que \eng{few} soit négatif, \eng{quite a few} exprime une quantité \textbf{importante}.
\end{definition}

\begin{exemplebox}[Exemples avec quite a few]
\begin{itemize}
\item \eng{Quite a few traders use Mobile Money in Douala.} « Pas mal de commerçants utilisent le Mobile Money à Douala. »
\item \eng{Quite a few students passed the English exam.} « Un bon nombre d'élèves ont réussi l'examen d'anglais. »
\item \eng{I have visited Buea quite a few times.} « J'ai visité Buéa pas mal de fois. »
\end{itemize}
\end{exemplebox}

\courssub{Few of / little of : devant the, my, these}

\begin{propriete}[La structure « of »]
Devant un déterminant (\eng{the}, \eng{my}, \eng{these}, \eng{those}) ou un pronom, on ajoute \cle{of} : \eng{a few of}, \eng{few of}, \eng{a little of}, \eng{little of}.
\begin{itemize}
\item \eng{A few of my classmates work on Saturdays.} « Quelques-uns de mes camarades travaillent le samedi. »
\item \eng{Few of the students understood the rule.} « Peu d'entre les élèves ont compris la règle. »
\item \eng{I drank a little of the juice.} « J'ai bu un peu du jus. »
\item \eng{A few of them sell phone credit.} « Quelques-uns d'entre eux vendent du crédit téléphonique. »
\end{itemize}
\end{propriete}

\courssub{Emploi pronominal : few et little sans nom}

\begin{propriete}[Few / a few / little / a little employés seuls]
Ces mots peuvent remplacer le nom entier quand le contexte est clair, notamment dans les réponses courtes.
\begin{itemize}
\item \eng{How many candidates came? — Few.} « Combien de candidats sont venus ? — Peu. »
\item \eng{Do you have friends in Bamenda? — A few.} « As-tu des amis à Bamenda ? — Quelques-uns. »
\item \eng{Is there any petrol left? — A little.} « Reste-t-il de l'essence ? — Un peu. »
\end{itemize}
\end{propriete}

\begin{savaistu}[Et « the few » ?]
Le groupe \eng{the few} existe aussi : \eng{the few books I own} signifie « les quelques livres que je possède ». Ici, \eng{the} rend le petit groupe précis et identifiable : \eng{The few customers who came this morning bought a lot.} « Les quelques clients venus ce matin ont beaucoup acheté. »
\end{savaistu}

\begin{exoresolu}[Choisir entre few, a few, little, a little]
Énoncé. Complétez avec \eng{few}, \eng{a few}, \eng{little} ou \eng{a little}, puis traduisez.
\begin{enumerate}
\item I have \ldots\ldots friends in Kumbo, so I never feel lonely there.
\item \ldots\ldots passengers were on the bus; it was almost empty.
\item Add \ldots\ldots salt to the sauce, please.
\item There is \ldots\ldots rain in the Far North this year; farmers are suffering.
\item \ldots\ldots of the pupils brought their dictionaries.
\end{enumerate}
\tcblower
Solution détaillée.
\begin{enumerate}
\item \eng{a few friends} : nom comptable pluriel, sens positif (« je ne me sens jamais seul »). « J'ai quelques amis à Kumbo, donc je ne m'y sens jamais seul. »
\item \eng{Few passengers} : nom comptable, sens négatif (« le bus était presque vide »). « Peu de passagers étaient dans le bus ; il était presque vide. »
\item \eng{a little salt} : nom incomptable, sens positif (on veut du sel). « Ajoute un peu de sel à la sauce, s'il te plaît. »
\item \eng{little rain} : nom incomptable, sens négatif (« les agriculteurs souffrent »). « Il y a peu de pluie dans l'Extrême-Nord cette année ; les agriculteurs souffrent. »
\item \eng{A few of the pupils} : structure avec \eng{of} devant \eng{the}. « Quelques-uns des élèves ont apporté leurs dictionnaires. »
\end{enumerate}
\end{exoresolu}

\begin{attention}[Pièges finaux]
\begin{itemize}
\item Ne confondez pas \eng{a little} (quantité, devant un nom) et \eng{a little} devant un adjectif : \eng{I am a little tired} « je suis un peu fatigué » — ici c'est correct, car il modifie l'adjectif, pas un nom comptable.
\item En français, « peu » est unique ; en anglais, choisissez toujours selon le \textbf{type de nom} (comptable ou non) et selon votre \textbf{intention} (positive ou négative).
\item \eng{Few} et \eng{little} sont déjà négatifs : ne doublez pas la négation. Faux : \eng{There isn't little hope.} Correct : \eng{There is little hope.}
\end{itemize}
\end{attention}

\begin{exercice}[À vous de jouer]
Traduisez en anglais : 1) Peu de clients sont venus aujourd'hui. 2) J'ai quelques amis qui parlent anglais. 3) Il reste un peu de temps avant la fermeture. 4) Pas mal de commerçants acceptent le Mobile Money. 5) Quelques-uns de mes voisins vendent des légumes au marché.
\end{exercice}

\coursec{All, both, each, every, most}

Dans la section précédente, nous avons étudié les quantifieurs qui expriment une \emph{petite quantité} : \eng{few}, \eng{a few}, \eng{little}, \eng{a little}. Nous passons maintenant aux déterminants qui expriment la \emph{totalité} ou une \emph{grande partie} d'un ensemble : \eng{all}, \eng{both}, \eng{each}, \eng{every} et \eng{most}. Ces mots sont extrêmement fréquents dans la vie économique et professionnelle : horaires d'ouverture, annonces commerciales, règlements d'entreprise. Observons d'abord un texte authentique.

\courssub{Observation et découverte guidée}

\begin{textebox}[A busy market day in Bamenda]
Every morning, all the traders of Nkwen Market arrive before seven o'clock. Most of them sell food: vegetables, plantains, beans and palm oil. Both of the big shops near the entrance open at half past seven, and each shopkeeper arranges his goods carefully. Every customer wants fresh products, so everybody comes early. “All my tomatoes are sold before noon,” says Mama Grace, a trader. “Most of my clients are regular customers. Each of them knows my prices.” Not all the stalls are busy at the same time: most of the activity happens in the morning. Both of her daughters help her at the weekend, and they all work hard. “Everything must be clean,” she adds. “Every plate, every basket. Each customer is important to me.” At the end of the day, most of the money is kept in a small metal box, and all of it is counted carefully before the family goes home. Everybody in the market knows Mama Grace, because she greets every person with a smile.
\end{textebox}

\textbf{Glossaire}

\begin{center}
\begin{tabularx}{\linewidth}{|X|X|X|}
\hline
\rowcolor{popblueL} Mot anglais & Nature & Traduction \\
\hline
trader & nom & commerçant(e) \\
\hline
stall & nom & étal, stand \\
\hline
shopkeeper & nom & boutiquier, commerçant \\
\hline
goods & nom pluriel & marchandises \\
\hline
to greet & verbe & saluer, accueillir \\
\hline
regular customer & expression & client habituel \\
\hline
\end{tabularx}
\end{center}

\textbf{Questions de compréhension}
\begin{enumerate}
\item What time do the traders arrive at Nkwen Market?
\item How many big shops are there near the entrance? How do you know?
\item When is most of the activity in the market?
\item Who helps Mama Grace at the weekend?
\end{enumerate}

Relevons maintenant les phrases-clés du texte et observons la place et la forme de chaque déterminant :

\begin{exemplebox}[Observation]
\begin{itemize}
\item \eng{All the traders arrive before seven.} — Tous les commerçants arrivent avant sept heures.
\item \eng{Both of the big shops open at half past seven.} — Les deux grandes boutiques ouvrent à sept heures et demie.
\item \eng{Each shopkeeper arranges his goods.} — Chaque commerçant range ses marchandises.
\item \eng{Every customer wants fresh products.} — Chaque client veut des produits frais.
\item \eng{Most of my clients are regular customers.} — La plupart de mes clients sont des habitués.
\item \eng{They all work hard.} — Ils travaillent tous dur.
\item \eng{Everybody comes early.} — Tout le monde vient tôt.
\end{itemize}
\end{exemplebox}

\courssub{All : la totalité}

\begin{propriete}[Formes et emplois de all]
\begin{itemize}
\item \cle{all} + nom \textbf{pluriel} ou \textbf{incomptable} (sans article) = « tout, tous, toute » en général : \eng{All students must work.} (Tous les élèves doivent travailler.) \eng{All water is precious.} (Toute l'eau est précieuse.)
\item \cle{all of} + \textbf{the / my / these / them...} (déterminant obligatoire après \eng{of}) : \eng{All of my friends are here.} (Tous mes amis sont là.) \eng{All of the money was counted.} (Tout l'argent a été compté.)
\item Avec un nom pluriel, le verbe est au \textbf{pluriel} ; avec un incomptable, au \textbf{singulier} : \eng{All the shops are closed.} / \eng{All the rice is sold.}
\end{itemize}
\end{propriete}

\begin{propriete}[Place de all quand il renvoie au sujet]
Quand \eng{all} renvoie au sujet, il se place \textbf{avant le verbe principal} mais \textbf{après le verbe be} et après l'auxiliaire :
\begin{itemize}
\item \eng{They all came to the meeting.} — Ils sont tous venus à la réunion.
\item \eng{They are all here.} — Ils sont tous là.
\item \eng{The workers have all been paid.} — Les ouvriers ont tous été payés (après l'auxiliaire).
\end{itemize}
\end{propriete}

\begin{attention}[All of students : erreur typique]
On ne dit jamais \eng{all of students}. Deux solutions correctes seulement : \eng{all students} (tous les élèves, en général) ou \eng{all of the students} (tous les élèves d'un groupe précis). De même : \eng{all children like games}, mais \eng{all of my children}.
\end{attention}

\courssub{Both : les deux}

\begin{propriete}[Both = exactement deux éléments]
\begin{itemize}
\item \cle{both} + nom \textbf{pluriel} : \eng{Both shops are closed.} (Les deux boutiques sont fermées.)
\item \cle{both of} + the / my / these / them : \eng{Both of my parents are teachers.} (Mes deux parents sont enseignants.) \eng{Both of them agreed.} (Tous les deux ont accepté.)
\item Le verbe est toujours au \textbf{pluriel}. \eng{Both} peut aussi suivre le sujet : \eng{They both live in Douala.} (Ils habitent tous les deux à Douala.)
\end{itemize}
\end{propriete}

\begin{exemplebox}[Exemples traduits]
\begin{itemize}
\item \eng{Both of my sisters are nurses.} — Mes deux sœurs sont infirmières.
\item \eng{Both banks open at eight o'clock.} — Les deux banques ouvrent à huit heures.
\item \eng{We both passed the examination.} — Nous avons tous les deux réussi l'examen.
\end{itemize}
\end{exemplebox}

\courssub{Each et every : chaque}

\begin{propriete}[Règle fondamentale de each et every]
\cle{each} et \cle{every} sont toujours suivis d'un nom au \textbf{SINGULIER} et d'un verbe au \textbf{singulier} :
\begin{itemize}
\item \eng{Each student has a book.} — Chaque élève a un livre.
\item \eng{Every shop opens at 8 a.m.} — Chaque boutique ouvre à 8 heures.
\end{itemize}
Différence de sens : \eng{each} considère les individus \textbf{un par un, séparément} (deux éléments ou plus) ; \eng{every} considère le \textbf{groupe dans son ensemble} (trois éléments ou plus), comme une généralisation.
\end{propriete}

\begin{propriete}[Each of et every one of]
\begin{itemize}
\item \cle{each of} + the / my / these + nom pluriel : \eng{Each of the students received a pen.} (Chacun des élèves a reçu un stylo.) Le verbe reste au singulier.
\item \eng{every} n'est \textbf{jamais} suivi de \eng{of}. On dit \cle{every one of} : \eng{Every one of them passed.} (Chacun d'entre eux a réussi.)
\end{itemize}
\end{propriete}

\begin{exemplebox}[Exemples traduits]
\begin{itemize}
\item \eng{Each pupil received a pen.} — Chaque élève a reçu un stylo.
\item \eng{Every worker has a contract.} — Chaque travailleur a un contrat.
\item \eng{Each of the farmers grows cocoa.} — Chacun des planteurs cultive le cacao.
\item \eng{Every one of the shops was open.} — Toutes les boutiques sans exception étaient ouvertes.
\end{itemize}
\end{exemplebox}

\begin{attention}[Every students : impossible]
L'erreur la plus fréquente des francophones est de mettre le nom au pluriel après \eng{every}, sous l'influence de « tous les ». On écrit \eng{every student}, \eng{every day}, \eng{every shop} — jamais \eng{every students}. Astuce : après \eng{every}, imaginez toujours « chaque », qui exige le singulier en français aussi.
\end{attention}

\begin{savaistu}[Each après le sujet]
\eng{Each} peut se placer après le sujet (ou après l'auxiliaire) ; le verbe redevient alors pluriel : \eng{They each received a gift.} = Ils ont chacun reçu un cadeau. \eng{The students each have a locker.} = Les élèves ont chacun un casier. C'est une structure très élégante à l'écrit.
\end{savaistu}

\courssub{Most : la plupart}

\begin{propriete}[Most et most of]
\begin{itemize}
\item \cle{most} + nom (sans article) = « la plupart de » en général : \eng{Most Cameroonians speak French.} (La plupart des Camerounais parlent français.) \eng{Most work is done by hand.} (La plupart du travail se fait à la main.)
\item \cle{most of} + the / my / these / them : \eng{Most of the students passed.} (La plupart des élèves ont réussi.) \eng{Most of the rice is sold.} (La plupart du riz est vendue.)
\item Jamais \eng{the most of} dans ce sens : \eng{the most} sert au superlatif (\eng{the most beautiful}).
\end{itemize}
\end{propriete}

\begin{exemplebox}[Exemples traduits]
\begin{itemize}
\item \eng{Most of the rice is sold.} — La plupart du riz est vendu.
\item \eng{Most traders accept mobile money.} — La plupart des commerçants acceptent le mobile money.
\item \eng{Most of them arrived on time.} — La plupart d'entre eux sont arrivés à l'heure.
\item \eng{Most of them work hard.} — La plupart d'entre eux travaillent dur.
\end{itemize}
\end{exemplebox}

\courssub{Les composés de every}

\begin{propriete}[Everybody, everyone, everything, everywhere]
Les composés de \eng{every} prennent un verbe au \textbf{SINGULIER} :
\begin{itemize}
\item \eng{Everybody is ready.} — Tout le monde est prêt.
\item \eng{Everyone has a right to work.} — Chacun a droit au travail.
\item \eng{Everything is expensive in town.} — Tout est cher en ville.
\item \eng{There are bendskins everywhere in Yaoundé.} — Il y a des bendskins partout à Yaoundé.
\end{itemize}
\end{propriete}

\begin{attention}[Everybody is, pas everybody are]
Sous l'influence de « tout le monde » perçu comme un groupe, les francophones écrivent souvent \eng{everybody are}. C'est faux : \eng{everybody is}, \eng{everyone knows}, \eng{everything costs}. Toujours le singulier.
\end{attention}

\begin{methode}[Comment choisir le bon déterminant de totalité ?]
\begin{enumerate}
\item \textbf{Nombre d'éléments :} Si c'est exactement \textbf{deux} $\rightarrow$ \eng{both}. Si c'est \textbf{trois ou plus} $\rightarrow$ \eng{all}, \eng{each}, \eng{every}, \eng{most}.
\item \textbf{Point de vue :} Individus séparés (un par un) $\rightarrow$ \eng{each}. Groupe entier (généralisation) $\rightarrow$ \eng{every}, \eng{all}. Grande majorité $\rightarrow$ \eng{most}.
\item \textbf{Nombre grammatical :} \eng{each}, \eng{every}, \eng{everybody/one/thing} $\rightarrow$ \textbf{singulier}. \eng{both}, \eng{all} (+ pluriel), \eng{most} (+ pluriel) $\rightarrow$ \textbf{pluriel}. \eng{all}, \eng{most} (+ incomptable) $\rightarrow$ \textbf{singulier}.
\item \textbf{Structure \eng{of} :} Après \eng{all of}, \eng{both of}, \eng{each of}, \eng{most of} $\rightarrow$ \textbf{déterminant obligatoire} (the, my, these, them...). Jamais \eng{every of} $\rightarrow$ \eng{every one of}.
\end{enumerate}
\end{methode}

\courssub{Tableau comparatif et synthèse}

\begin{popfigure}
\begin{tikzpicture}[
box/.style={draw=popline, rounded corners=3pt, fill=popcream, align=center, minimum width=3.4cm, minimum height=2.6cm, font=\small},
titre/.style={font=\bfseries\small, text=popdark}
]
\node[box, fill=popblueL, align=center] (each) at (0,0){\textbf{each}\\[2pt] 2 éléments ou +\\ un par un\\ nom singulier\\ \emph{each student}};
\node[box, fill=popgreenL, align=center] (every) at (4.4,0){\textbf{every}\\[2pt] 3 éléments ou +\\ le groupe entier\\ nom singulier\\ \emph{every shop}};
\node[box, fill=poporangeL, align=center] (all) at (8.8,0){\textbf{all}\\[2pt] la totalité\\ pluriel ou\\ incomptable\\ \emph{all students / all water}};
\node[box, fill=poppinkL, align=center] (both) at (13.2,0){\textbf{both}\\[2pt] exactement 2\\ nom pluriel\\ verbe pluriel\\ \emph{both shops}};
\node[titre] at (6.6,2.2) {Choisir le bon déterminant de totalité};
\end{tikzpicture}
\legende{Each, every, all, both : nombre d'éléments, point de vue et nombre grammatical.}
\end{popfigure}

\begin{center}
\begin{tabularx}{\linewidth}{|l|X|X|X|}
\hline
\rowcolor{popblueL} Déterminant & Nom qui suit & Verbe & Traduction \\
\hline
all (+ of the) & pluriel / incomptable & pluriel / singulier & tous, tout \\
\hline
both (of) & pluriel (2 éléments) & pluriel & les deux \\
\hline
each (of) & singulier & singulier & chaque (un par un) \\
\hline
every & singulier & singulier & chaque, tous les \\
\hline
most (of) & pluriel / incomptable & selon le nom & la plupart de \\
\hline
\end{tabularx}
\end{center}

\begin{center}
\begin{tabularx}{\linewidth}{|X|X|X|}
\hline
\rowcolor{popblueL} Français & English & Remarque \\
\hline
Tous les élèves & all students / all of the students & jamais \eng{all of students} \\
\hline
Chaque élève & each student / every student & nom au singulier \\
\hline
Les deux boutiques & both shops / both of the shops & verbe au pluriel \\
\hline
La plupart des clients & most customers / most of the customers & jamais \eng{the most of} \\
\hline
Tout le monde est là & everybody is here & verbe au singulier \\
\hline
\end{tabularx}
\end{center}

\begin{exoresolu}[Exercice résolu : complétez avec all, both, each, every, most]
Énoncé — Complétez les phrases avec le déterminant qui convient :
\begin{enumerate}
\item \eng{... employee must wear a uniform. (règle générale, un par un)}
\item \eng{... of my two brothers work in a bank.}
\item \eng{... the cocoa is exported from Douala. (totalité)}
\item \eng{... of the traders in this market are women. (grande partie)}
\item \eng{... of the thirty students has a dictionary.}
\end{enumerate}
\tcblower
Solution détaillée :
\begin{enumerate}
\item \eng{Every employee must wear a uniform.} — Règle générale visant le groupe : \eng{every} + singulier (\eng{each} serait aussi possible si l'on insiste sur l'individu).
\item \eng{Both of my brothers work in a bank.} — Exactement deux personnes, verbe au pluriel.
\item \eng{All the cocoa is exported from Douala.} — Nom incomptable, totalité, verbe au singulier.
\item \eng{Most of the traders in this market are women.} — Grande partie d'un groupe précis : \eng{most of the} + pluriel.
\item \eng{Each of the thirty students has a dictionary.} — \eng{each of} + pluriel, mais verbe au singulier (\eng{has}). \eng{Every of} serait impossible.
\end{enumerate}
\end{exoresolu}

\begin{exercice}[À vous de jouer]
Traduisez en anglais :
\begin{enumerate}
\item Mes deux sœurs sont infirmières.
\item Chaque élève a reçu un stylo.
\item La plupart du riz est vendu.
\item Ils travaillent tous dur.
\item Tout le monde est prêt.
\item Toutes les boutiques ouvrent à huit heures.
\end{enumerate}
\end{exercice}

\begin{corrige}[Exercice « À vous de jouer »]
\begin{enumerate}
\item \eng{Both of my sisters are nurses.}
\item \eng{Each pupil received a pen.}
\item \eng{Most of the rice is sold.}
\item \eng{They all work hard.}
\item \eng{Everybody is ready.}
\item \eng{Every shop opens at eight o'clock.} (ou \eng{All the shops open at eight o'clock.})
\end{enumerate}
\end{corrige}

\begin{aretenir}[L'essentiel]
\eng{all} = totalité (pluriel ou incomptable) ; \eng{both} = exactement deux (pluriel) ; \eng{each} et \eng{every} = chaque, toujours avec un nom et un verbe au singulier ; \eng{most} = la plupart. Après \eng{of}, un déterminant est obligatoire : \eng{all of the}, \eng{each of my}, \eng{most of them}. Les composés de \eng{every} prennent le singulier : \eng{everybody is}.
\end{aretenir}

\coursec{Synthèse, comparaison avec le français et entraînement}

Dans les sections précédentes, nous avons étudié en détail les quantifieurs anglais : \eng{some} et \eng{any}, \eng{no} et \eng{none}, \eng{much} et \eng{many}, \eng{few}, \eng{a few}, \eng{little}, \eng{a little}, puis \eng{all}, \eng{both}, \eng{each}, \eng{every} et \eng{most}. Il est temps maintenant de rassembler toutes ces pièces du puzzle dans un tableau unique, de comparer systématiquement le fonctionnement de l'anglais avec celui du français, puis de nous entraîner jusqu'à ce que le choix du bon quantifieur devienne un réflexe. C'est précisément l'objectif de cette dernière section : passer de la connaissance des règles à l'automatisme, comme l'exigent les épreuves de grammaire et de composition de la classe de Première.

\courssub{Le tableau récapitulatif général}

Avant tout, observons ce court dialogue entre deux élèves de première qui préparent une excursion au marché de Mokolo, à Yaoundé. Lis-le une première fois pour le sens général, puis une deuxième fois en soulignant tous les quantifieurs.

\begin{textebox}[At Mokolo Market — a dialogue]
\eng{— How much money do you have?\\
— Not much. Only a few thousand francs. But we don't need a lot of money: most of the food at Mokolo is cheap.\\
— Are there many customers today?\\
— Yes, there are a lot of people, but there are no queues at the tomato sellers. Each seller has a little table, and every table is full of vegetables.\\
— Good. I want to buy some plantains and a little palm oil.\\
— All right. Both of us can carry the bags. None of the bags will be too heavy!}
\end{textebox}

Glossaire : \emph{queue} (n., file d'attente) ; \emph{seller} (n., vendeur, vendeuse) ; \emph{plantain} (n., plantain) ; \emph{palm oil} (n., huile de palme) ; \emph{to carry} (v., porter, transporter) ; \emph{heavy} (adj., lourd).

Observons : chaque quantifieur du dialogue obéit à trois questions simples. C'est la règle d'or de cette leçon.

\begin{propriete}[La règle d'or : trois questions pour choisir le bon quantifieur]
Avant de choisir un quantifieur, pose-toi toujours trois questions, dans cet ordre :
\begin{enumerate}
\item \cle{Comptable ou incomptable ?} Le nom peut-il se compter (\eng{customer}, \eng{bag}) ou non (\eng{money}, \eng{oil}) ?
\item \cle{Affirmative, négative ou interrogative ?} La phrase est-elle affirmative, négative ou interrogative ?
\item \cle{Sens positif ou négatif ?} Veut-on dire « il y en a » (\eng{a few}, \eng{a little}) ou « il n'y en a presque pas » (\eng{few}, \eng{little}) ?
\end{enumerate}
Une fois ces trois réponses trouvées, le bon quantifieur s'impose automatiquement.
\end{propriete}

\begin{popfigure}
\begin{tikzpicture}[font=\small]
\draw[fill=popblueL, draw=popblue] (0,0) rectangle (2.6,0.9);
\node[align=center] at (1.3,0.45) {\textbf{Quantifier}};
\draw[fill=poporangeL, draw=poporange] (2.6,0) rectangle (5.6,0.9);
\node[align=center] at (4.1,0.45) {\textbf{Noun type}};
\draw[fill=popgreenL, draw=popgreen] (5.6,0) rectangle (8.6,0.9);
\node[align=center] at (7.1,0.45) {\textbf{Sentence type}};
\draw[fill=poppurpleL, draw=poppurple] (8.6,0) rectangle (12.6,0.9);
\node[align=center] at (10.6,0.45) {\textbf{Example}};
\foreach \y/\a/\b/\c/\d in {
-0.7/{some}/{countable + uncountable}/{affirmative (+ offers)}/{I bought some plantains.},
-1.4/{any}/{countable + uncountable}/{negative + interrogative}/{There isn't any oil.},
-2.1/{no, none}/{countable + uncountable}/{affirmative form; negative meaning}/{There are no customers. None left.},
-2.8/{much}/{uncountable}/{negative + interrogative}/{How much money do you have?},
-3.5/{many}/{countable plural}/{all types}/{Many people were at the market.},
-4.2/{a few, few}/{countable plural}/{positive vs negative meaning}/{A few sellers smiled. Few waited.},
-4.9/{a little, little}/{uncountable}/{positive vs negative meaning}/{A little oil. Little time.},
-5.6/{each, every}/{singular countable}/{affirmative}/{Each table is full. Every day.},
-6.3/{all, both}/{countable + uncountable}/{affirmative}/{All the food. Both of us.},
-7.0/{most (of)}/{countable + uncountable}/{affirmative}/{Most of the food is cheap.}}
{
\draw[draw=popline] (0,\y) rectangle (2.6,\y+0.7); \node[align=center] at (1.3,\y+0.35) {\a};
\draw[draw=popline] (2.6,\y) rectangle (5.6,\y+0.7); \node[align=center] at (4.1,\y+0.35) {\b};
\draw[draw=popline] (5.6,\y) rectangle (8.6,\y+0.7); \node[align=center] at (7.1,\y+0.35) {\c};
\draw[draw=popline] (8.6,\y) rectangle (12.6,\y+0.7); \node[anchor=west, text width=3.7cm, align=flush left, inner sep=1pt] at (8.7,\y+0.35) {\d};
}
\end{tikzpicture}
\legende{Tableau de synthèse : le choix du quantifieur dépend du type de nom et du type de phrase.}
\end{popfigure}

\begin{aretenir}[Of ou pas of ? La question qui change tout]
Devant un nom, le quantifieur s'emploie \emph{sans} \eng{of} : \eng{some water}, \eng{many students}, \eng{most people}. Mais devant \eng{the}, \eng{my}, \eng{these} ou un pronom, il faut \eng{of} : \eng{some of the water}, \eng{many of my students}, \eng{most of them}, \eng{both of us}, \eng{none of the bags}. Compare : \eng{Most students passed.} « La plupart des élèves ont réussi » (en général) et \eng{Most of the students passed.} « La plupart des élèves (de cette classe) ont réussi » (groupe précis).
\end{aretenir}

\courssub{Comparaison systématique avec le français}

Le français et l'anglais ne découpent pas la réalité de la même façon. Voici les quatre grandes différences à maîtriser absolument.

\begin{propriete}[L'article partitif français n'existe pas en anglais]
Le français utilise l'article partitif (\emph{du}, \emph{de la}, \emph{des}) là où l'anglais utilise \cle{some} ou l'\cle{article zéro} (c'est-à-dire rien du tout).
\begin{itemize}
\item \eng{I drink water.} « Je bois de l'eau. » (vérité générale : article zéro)
\item \eng{I drink some water every morning.} « Je bois de l'eau chaque matin. » (quantité indéfinie : \eng{some})
\item \eng{Do you want some tea?} « Veux-tu du thé ? »
\end{itemize}
En anglais, on ne met jamais \emph{the} pour exprimer le partitif : \eng{I drink the water} signifie « je bois \emph{l'}eau » (une eau précise, déjà identifiée).
\end{propriete}

\begin{propriete}[« Beaucoup de » : much, many ou a lot of]
« Beaucoup de » se traduit de trois façons selon le nom et le type de phrase :
\begin{itemize}
\item \cle{many} + nom comptable pluriel : \eng{many students} « beaucoup d'élèves » ;
\item \cle{much} + nom incomptable, surtout à la négative et à l'interrogative : \eng{There isn't much time.} « Il ne reste pas beaucoup de temps. » ;
\item \cle{a lot of} (ou \eng{lots of}) + les deux types de noms, surtout à l'affirmative : \eng{a lot of money}, \eng{a lot of people}.
\end{itemize}
On ne dit jamais « much people » : \eng{much} ne s'emploie pas avec un nom comptable pluriel.
\end{propriete}

\begin{propriete}[Each, every, all, both : l'accord du verbe]
\begin{itemize}
\item \cle{each} et \cle{every} + nom \emph{singulier} + verbe au \emph{singulier} : \eng{Every student has a book.} « Chaque élève a un livre. » \eng{Each of the students has a book.} « Chacun des élèves a un livre. »
\item \cle{all} + nom pluriel ou incomptable : \eng{All the students are here.} « Tous les élèves sont là. » \eng{All the food is cheap.} « Toute la nourriture est bon marché. »
\item \cle{both} = exactement deux, toujours suivi d'un pluriel : \eng{Both brothers are traders.} « Les deux frères sont commerçants. »
\item \cle{none of} + nom pluriel : le verbe peut être au singulier (style soigné) ou au pluriel (style courant) : \eng{None of the shops is open.} ou \eng{None of the shops are open.} « Aucune des boutiques n'est ouverte. »
\end{itemize}
\end{propriete}

\begin{attention}[Les quatre pièges majeurs pour les francophones]
\begin{enumerate}
\item Ne traduis jamais « beaucoup d'argent » par « many money » : \eng{money} est incomptable. Dis \eng{a lot of money} ou \eng{much money}.
\item « La plupart des élèves » = \eng{most of the students} ou \eng{most students}. Jamais « the most of the students » : \eng{the most} sert au superlatif (\eng{the most beautiful}).
\item « Aucun » seul (pronom) = \cle{none} : \eng{How many pens do you have? — None.} « Aucun + nom » = \cle{no} + nom : \eng{I have no pens.} On n'écrit jamais « none pens ».
\item « Chaque jour » = \eng{every day} (en deux mots) ; \eng{everyday} (en un mot) est un adjectif : \eng{my everyday life} « ma vie quotidienne ».
\end{enumerate}
\end{attention}

\begin{center}
\begin{tabularx}{\linewidth}{|X|X|X|}
\hline
\rowcolor{popblueL} Français & English & Remarque \\
\hline
Je bois de l'eau. & I drink water. / I drink some water. & article zéro ou \eng{some}, jamais \eng{the} \\
\hline
Beaucoup d'argent & a lot of money / much money & jamais « many money » \\
\hline
Beaucoup de clients & many customers / a lot of customers & nom comptable pluriel \\
\hline
La plupart des élèves & most of the students / most students & jamais « the most of » \\
\hline
Il n'y a aucun client. & There are no customers. / There aren't any customers. & deux formulations possibles \\
\hline
Aucun. (réponse seule) & None. & pronom, sans nom derrière \\
\hline
Chaque jour & every day & deux mots ; suivi d'un verbe au singulier \\
\hline
Les deux frères & both brothers / both of the brothers & \eng{both} = exactement deux \\
\hline
Quelques amis / peu d'amis & a few friends / few friends & sens positif / sens négatif \\
\hline
Un peu de temps / peu de temps & a little time / little time & sens positif / sens négatif \\
\hline
\end{tabularx}
\end{center}

\begin{exemplebox}[Exemples traduits à mémoriser]
\begin{itemize}
\item \eng{There are no customers today.} « Il n'y a aucun client aujourd'hui. »
\item \eng{There aren't any customers today.} (même sens, avec \eng{any})
\item \eng{Most of the students passed the exam.} « La plupart des élèves ont réussi l'examen. »
\item \eng{Every day, she sells a few tomatoes and a little oil.} « Chaque jour, elle vend quelques tomates et un peu d'huile. »
\item \eng{Both of my uncles are traders in Douala.} « Mes deux oncles sont commerçants à Douala. »
\item \eng{None of the shops is open on Sunday.} « Aucune des boutiques n'est ouverte le dimanche. »
\item \eng{Few passengers had tickets, but a few had passes.} « Peu de passagers avaient des tickets, mais quelques-uns avaient des laissez-passer. »
\end{itemize}
\end{exemplebox}

\begin{savaistu}[« A lot of » à l'oral]
À l'oral, les anglophones préfèrent \eng{a lot of} à \eng{much} dans les phrases affirmatives : \eng{I have a lot of work} sonne beaucoup plus naturel que \eng{I have much work}, qui fait très écrit, voire vieilli. À l'inverse, à la négative et à l'interrogative, \eng{much} et \eng{many} restent parfaitement normaux : \eng{Do you have much homework?} Autre détail sonore : \eng{a lot of} se prononce souvent « a-lo-teuv » à vitesse normale, comme un seul mot.
\end{savaistu}

\courssub{Méthode : réussir les exercices de transformation}

\begin{methode}[Transformer une phrase affirmative en négative ou interrogative]
\begin{enumerate}
\item Repère le quantifieur de la phrase affirmative : s'il s'agit de \eng{some}, il devient presque toujours \eng{any} à la négative et à l'interrogative.
\item Vérifie le type de nom : comptable pluriel (\eng{many}, \eng{a few}) ou incomptable (\eng{much}, \eng{a little}) ? Le quantifieur doit rester compatible.
\item Choisis la formulation négative : soit \eng{not ... any}, soit \eng{no} + nom (une seule négation par phrase, jamais les deux).
\item À l'interrogative, garde \eng{some} uniquement pour une offre ou une demande polie : \eng{Would you like some tea?}
\item Relis la phrase obtenue et vérifie l'accord du verbe : \eng{every} + singulier, \eng{few people} + pluriel, \eng{much time} + singulier.
\end{enumerate}
Exemple : \eng{There is some rice.} → négative : \eng{There isn't any rice.} ou \eng{There is no rice.} → interrogative : \eng{Is there any rice?}
\end{methode}

\begin{exoresolu}[Exercice résolu : choisir le bon quantifieur]
Énoncé — Complète avec \eng{much}, \eng{many}, \eng{a few}, \eng{a little}, \eng{no}, \eng{none}, \eng{most of the} ou \eng{every} :\\
1. How \ldots\space money do you earn per month? 2. \ldots\space trader in this market pays a small tax. 3. I have \ldots\space friends in Bamenda; we meet every week. 4. There is \ldots\space bread left; we must buy some. 5. \ldots\space students in my class speak English well. 6. « How many chairs are free? » — « \ldots . They are all occupied. »
\tcblower
Solution détaillée :\\
1. \eng{much} : \eng{money} est incomptable et la phrase est interrogative.\\
2. \eng{Every} : nom singulier (\eng{trader}), sens de « chaque » ; le verbe \eng{pays} est au singulier, ce qui confirme.\\
3. \eng{a few} : \eng{friends} est comptable pluriel et le sens est positif (« j'en ai quelques-uns »).\\
4. \eng{no} : « il ne reste pas de pain » ; on peut aussi dire \eng{There isn't any bread left}.\\
5. \eng{Most of the} : « la plupart des élèves » ; jamais « the most of ».\\
6. \eng{None} : réponse seule, sans nom après ; « aucune ». « No » serait impossible ici car \eng{no} exige un nom.
\end{exoresolu}

\courssub{Entraînement}

\begin{exercice}[Exercice 1 — À trous]
Complète avec le quantifieur correct : \eng{some}, \eng{any}, \eng{no}, \eng{much}, \eng{many}, \eng{a few}, \eng{a little}, \eng{each}, \eng{most}.\\
1. There isn't \ldots\space sugar in my tea. 2. \ldots\space of the farmers here grow cocoa. 3. She speaks \ldots\space French, but not fluently. 4. How \ldots\space books did you buy? 5. Would you like \ldots\space coffee? 6. \ldots\space student received a certificate. 7. We have \ldots\space time; the bus leaves in two minutes! 8. Only \ldots\space passengers were on the bendskin.
\end{exercice}

\begin{exercice}[Exercice 2 — Choix multiples]
Entoure la bonne réponse.\\
1. I don't have (much / many / a few) money. 2. (The most of / Most of the / Most) shops in this street close at six. 3. There are (no / none / not) tickets left. 4. « How many seats are free? » — « (No / None / Any). » 5. (Every / All / Both) day, my father goes to the market. 6. She drank (a few / a little / many) water.
\end{exercice}

\begin{exercice}[Exercice 3 — Transformation]
Mets les phrases à la forme négative, puis à la forme interrogative.\\
1. There is some milk in the fridge. 2. She has many customers. 3. He earns much money. 4. There are some taxis at the station.
\end{exercice}

\begin{exercice}[Exercice 4 — Traduction]
Traduis en anglais : 1. Je bois de l'eau tous les jours. 2. Il n'y a aucun client dans la boutique. 3. La plupart des élèves aiment le football. 4. Nous avons peu de temps mais un peu d'argent. 5. Chaque vendeur a quelques tomates.
\end{exercice}

\begin{exercice}[Exercice 5 — Production écrite : « My local market »]
Rédige un paragraphe de huit à dix phrases intitulé \eng{My local market}. Utilise au moins huit quantifieurs différents parmi : \eng{some}, \eng{any}, \eng{no}, \eng{none}, \eng{much}, \eng{many}, \eng{a few}, \eng{a little}, \eng{each}, \eng{every}, \eng{all}, \eng{both}, \eng{most of the}. Souligne chaque quantifieur utilisé. Décris les vendeurs, les clients, les produits et les prix.
\end{exercice}

\begin{corrige}[Exercices 1 à 4]
\textbf{Exercice 1 :} 1. \eng{any} 2. \eng{Most} 3. \eng{a little} 4. \eng{many} 5. \eng{some} (offre polie) 6. \eng{Each} 7. \eng{no} 8. \eng{a few}.\\
\textbf{Exercice 2 :} 1. \eng{much} 2. \eng{Most of the} (groupe précis : « les boutiques de cette rue ») 3. \eng{no} 4. \eng{None} 5. \eng{Every} 6. \eng{a little}.\\
\textbf{Exercice 3 :} 1. \eng{There isn't any milk in the fridge. / There is no milk in the fridge. Is there any milk in the fridge?} 2. \eng{She doesn't have many customers. Does she have many customers?} 3. \eng{He doesn't earn much money. Does he earn much money?} 4. \eng{There aren't any taxis at the station. / There are no taxis at the station. Are there any taxis at the station?}\\
\textbf{Exercice 4 :} 1. \eng{I drink water every day.} 2. \eng{There are no customers in the shop. / There aren't any customers in the shop.} 3. \eng{Most of the students like football.} 4. \eng{We have little time but a little money.} 5. \eng{Each seller has a few tomatoes.}
\end{corrige}

\begin{aretenir}[L'essentiel de la leçon]
Retiens la règle d'or : \cle{comptable ou incomptable ? affirmative, négative ou interrogative ? sens positif ou négatif ?} Trois réponses, et le quantifieur s'impose de lui-même. Garde en tête les quatre pièges : jamais « many money », jamais « the most of », \eng{none} seul mais \eng{no} + nom, et le partitif français qui devient \eng{some} ou l'article zéro. N'oublie pas non plus la règle du \eng{of} : \eng{most people} mais \eng{most of the people}, \eng{both brothers} mais \eng{both of us}. Avec le tableau de synthèse sous les yeux et un peu d'entraînement quotidien, les quantifieurs anglais n'auront bientôt plus aucun secret pour toi.
\end{aretenir}

\newpage

\coursec{Activité d'intégration}

\courssub{Objectifs de l'activité}

À l'issue de cette activité d'intégration, l'élève sera capable de :
\begin{itemize}
\item utiliser spontanément, à l'oral, au moins \cle{huit quantifieurs différents} (\eng{some, any, no, none, much, many, a few, a little, most, both, all, each, every}) dans un échange authentique de type achat au marché ;
\item poser des questions correctes avec \eng{how much} (devant un nom indénombrable) et \eng{how many} (devant un nom dénombrable au pluriel) ;
\item rédiger un court compte rendu écrit (8 à 10 lignes) en réemployant ces quantifieurs dans un texte cohérent au prétérit ;
\item évaluer la production de ses camarades à l'aide d'une grille d'observation simple (compétence d'\emph{auto-} et \emph{co-évaluation}).
\end{itemize}

\courssub{Situation}

It is Saturday morning. Your mother has given you 5,000 FCFA and a shopping list. You are at Mokolo Market in Yaoundé, one of the busiest markets in the city. You must buy food for the family, ask about prices, negotiate with the trader, and check that you have enough money. Work in groups of three: one \eng{customer}, one \eng{trader}, one \eng{observer}.

\begin{textebox}[The customer's shopping list and budget]
\textbf{Shopping list — Saturday, Mokolo Market}
\begin{itemize}
\item rice (2 kg) --- groundnut oil (1 litre)
\item tomatoes (1 kg) --- onions (half a kilo)
\item bananas (a bunch) --- soap (2 bars)
\item okra (a small basket)
\end{itemize}
\textbf{Budget:} 5,000 FCFA maximum. Do not spend all the money: keep some for the \emph{bendskin} (motorbike taxi) back home!
\end{textebox}

\begin{textebox}[The trader's stall — what is available today]
\textbf{Mama Ngo Bell's stall}
\begin{itemize}
\item Rice: plenty (700 FCFA/kg) --- Oil: only a little left (1,200 FCFA/litre)
\item Tomatoes: a few, very fresh (500 FCFA/kg) --- Onions: none left, sorry!
\item Bananas: many, both sweet bananas and plantains (300 FCFA/bunch)
\item Soap: some bars (250 FCFA each) --- Okra: no okra today, come back on Tuesday.
\end{itemize}
\textbf{Trader's tip:} Most customers pay with Mobile Money, but cash is fine too.
\end{textebox}

\courssub{Consignes}

\begin{enumerate}
\item \textbf{Préparation (5 min).} Dans votre groupe, distribuez les rôles : \eng{customer}, \eng{trader}, \eng{observer}. Le client relit sa liste et son budget ; le vendeur relit l'état de son stock ; l'observateur prépare sa grille (voir ci-dessous).
\item \textbf{Jeu de rôle (7 min).} Jouez la scène en anglais. Le client salue (\eng{Good morning, Madam!}), demande les prix avec \eng{how much} ou \eng{how many}, négocie (\eng{That's too much! Can you reduce the price a little?}), achète ou renonce. Le vendeur répond en utilisant les quantifieurs : \eng{We have no onions left, but there are a few tomatoes. Most customers pay with Mobile Money. Both soaps are good quality.} L'observateur coche chaque quantifieur correctement employé et note les erreurs.
\item \textbf{Grille de l'observateur.} Cochez à chaque emploi correct :

\begin{center}
\begin{tabular}{|l|c|l|c|}
\hline
\rowcolor{popblueL} \textbf{Quantifier} & \textbf{Used correctly?} & \textbf{Quantifier} & \textbf{Used correctly?} \\
\hline
some & $\square$ & no / none & $\square$ \\
\hline
any & $\square$ & a few / few & $\square$ \\
\hline
much & $\square$ & a little / little & $\square$ \\
\hline
many & $\square$ & most & $\square$ \\
\hline
how much & $\square$ & both / all & $\square$ \\
\hline
how many & $\square$ & each / every & $\square$ \\
\hline
\end{tabular}
\end{center}

\item \textbf{Compte rendu écrit (10 min).} Individuellement ou en binôme, rédigez un texte de 8 à 10 lignes intitulé \eng{Our shopping at the market}. Racontez la scène au prétérit et utilisez \textbf{au moins huit quantifieurs différents}. Soulignez-les dans votre copie.
\item \textbf{Restitution (5 min).} Le groupe qui a utilisé le plus de quantifieurs corrects joue son dialogue devant la classe. La classe vérifie avec la grille.
\end{enumerate}

\courssub{Ressources à mobiliser}

\begin{aretenir}[Rappel express des règles]
\begin{itemize}
\item \eng{some} : phrase affirmative (\eng{I bought some rice.}) ; \eng{any} : phrase négative et interrogative (\eng{Do you have any soap? We don't have any onions.}). \textbf{Exception :} on utilise \eng{some} dans les questions qui sont des offres ou des demandes polies : \eng{Would you like some tea? Can I have some change, please?}
\item \eng{no} + nom = \eng{not... any} (\eng{We have no onions left.} = \eng{We don't have any onions left.}) ; \eng{none} : pronom, employé seul, sans nom derrière (\eng{How many onions are left? None.}).
\item \eng{much} / \eng{how much} + nom indénombrable (\eng{How much is the oil?}) ; \eng{many} / \eng{how many} + nom dénombrable au pluriel (\eng{How many bananas do you want?}).
\item \eng{a few} / \eng{a little} = « quelques / un peu » (sens positif : il y en a assez) ; \eng{few} / \eng{little} = « peu » (sens négatif : presque pas).
\item \eng{most} + nom pluriel ou indénombrable, sans article (\eng{Most customers pay cash.}) ; \eng{both} = « les deux » (+ pluriel) ; \eng{each} + nom singulier et verbe au singulier ; \eng{every} + nom singulier.
\end{itemize}
\end{aretenir}

\textbf{Lexique utile du marché :} \eng{a stall} (un étal), \eng{a bunch of bananas} (un régime de bananes), \eng{a bar of soap} (une savonnette), \eng{fresh} (frais, fraîche), \eng{to bargain / to negotiate} (marchander), \eng{change} (la monnaie — indénombrable), \eng{left} (restant : \eng{no onions left} = « plus d'oignons »).

\begin{attention}[Pièges à éviter]
Ne dites pas \eng{How much tomatoes?} mais \eng{How many tomatoes?} (nom dénombrable au pluriel). Ne dites pas \eng{I don't have some money} mais \eng{I don't have any money}. Après \eng{no}, pas de négation supplémentaire : \eng{We have no onions}, jamais \eng{We don't have no onions} (la double négation est fautive en anglais standard). \eng{Money} est indénombrable : \eng{too much money}, jamais \eng{too many money}. Enfin, \eng{each} exige un verbe au singulier : \eng{Each bar costs 250 FCFA.}
\end{attention}

\courssub{Proposition de production et corrigé commenté}

\begin{exoresolu}[Our shopping at the market — modèle de compte rendu]
Rédigez un compte rendu de 8 à 10 lignes utilisant au moins huit quantifieurs différents.
\tcblower
\textbf{Model answer:}

\eng{Last Saturday, we went shopping at Mokolo Market with a budget of 5,000 FCFA. There were many people, and most stalls were already busy. First, we bought some rice and a few tomatoes, because they were very fresh. Unfortunately, the trader had no onions left, so we could not buy any. We also asked how much the groundnut oil cost, but there was only a little left, so we took one litre. Both bars of soap were cheap, and each cost 250 FCFA. We did not spend all our money: we kept some for the bendskin home. Every member of the family was happy with our shopping!}

\textbf{Commentaire des choix de langue :}
\begin{itemize}
\item \eng{many people}, \eng{most stalls} : noms dénombrables au pluriel — correct ;
\item \eng{some rice}, \eng{a few tomatoes} : indénombrable avec \eng{some}, dénombrable pluriel avec \eng{a few} (sens positif) ;
\item \eng{no onions left} puis \eng{any} dans la phrase négative (\eng{could not buy any}) : les deux formes de la négation sont illustrées ;
\item \eng{how much the groundnut oil cost} : question indirecte — attention à l'ordre des mots (sujet + verbe, sans inversion) ;
\item \eng{a little left} : indénombrable, sens positif (« il en restait un peu ») ;
\item \eng{both bars}, \eng{each cost} : \eng{both} + pluriel ; \eng{each} + verbe au singulier (\eng{cost} est ici le prétérit de \eng{cost}, identique au présent) ;
\item \eng{all our money}, \eng{some} (pronom, seul), \eng{every member} + singulier : au total, douze quantifieurs différents — objectif dépassé.
\end{itemize}
\end{exoresolu}

\courssub{Bilan de l'activité}

Cette activité a permis de réinvestir l'ensemble des quantifieurs de la leçon dans une situation de communication réelle et familière : le marché. À l'oral, vous avez appris à choisir entre \eng{how much} et \eng{how many} selon la nature du nom (indénombrable ou dénombrable), et à employer \eng{no}, \eng{none}, \eng{a few}, \eng{a little} pour décrire un stock. À l'écrit, vous avez mobilisé ces mêmes outils au prétérit dans un texte cohérent, compétence directement réinvestissable dans l'épreuve de \emph{composition} de l'examen. La grille de l'observateur vous a rendus acteurs de votre évaluation. Critère de réussite : au moins huit quantifieurs différents, correctement employés, dans le dialogue et dans le compte rendu. Si ce critère n'est pas atteint, reprenez les tableaux de la leçon et refaites le jeu de rôle en changeant les rôles.

\newpage

\coursec{Méthodes et exercices résolus}

\courssub{Méthode 1 : Identifier le déterminant correct}

\begin{methode}[Identifier et construire correctement la structure étudiée]
Pour choisir le bon déterminant ou quantifieur, suivre une démarche en quatre étapes :

\begin{enumerate}
\item \cle{Identifier le nom} qui suit : est-il \textbf{comptable} (\eng{books, students, markets}) ou \textbf{incomptable} (\eng{money, water, information}) ?
\item \cle{Identifier le type de phrase} : affirmative, négative ou interrogative ?
\item \cle{Identifier le sens voulu} : quantité positive (\eng{a few / a little} = « quelques, un peu de »), quantité quasi nulle (\eng{few / little} = « peu de »), totalité (\eng{all, every}), dualité (\eng{both}), absence (\eng{no, none}) ?
\item \cle{Vérifier l'accord} : le verbe s'accorde-t-il au singulier (\eng{each, every, little, much}) ou au pluriel (\eng{many, few, both, all} + nom pluriel) ?
\end{enumerate}

\textbf{Réflexes à retenir :}
\begin{itemize}
\item \eng{much} + incomptable ; \eng{many} + pluriel comptable.
\item \eng{some} en affirmatif, \eng{any} en négatif et interrogatif (sauf offre ou demande polie : \eng{Would you like some tea?} « Voulez-vous du thé ? »).
\item \eng{no} + nom (sans négation du verbe) ; \eng{none} seul ou \eng{none of} + déterminant + nom.
\item \eng{each} et \eng{every} + nom singulier + verbe singulier.
\item \eng{of} est obligatoire devant un déterminant (\eng{the, my, these}) ou un pronom : \eng{most of the students}, \eng{both of them} ; il disparaît sans déterminant : \eng{most students}, \eng{both brothers}.
\end{itemize}
\end{methode}

\begin{exoresolu}[Exercice type : choix du quantifieur]
\textbf{Énoncé.} Complete each sentence with the correct determiner: \emph{much, many, few, little, some, any, no, each}.

\begin{enumerate}
\item There isn't \dots{} rice left in the bag; go to the market in Mokolo.
\item \dots{} trader in the market pays a daily tax.
\item How \dots{} students passed the English test?
\item We have very \dots{} time before the bus leaves for Buea.
\item She earns \dots{} money selling vegetables, but she is happy.
\item There are \dots{} chairs in the classroom; we need ten more.
\end{enumerate}

\tcblower
\textbf{Solution détaillée.}

\begin{enumerate}
\item \eng{There isn't \textbf{any} rice left.} — Phrase négative (\eng{isn't}) + nom incomptable (\eng{rice}) : on emploie \cle{any}. \eng{Much} serait possible (\eng{There isn't much rice left}) avec une nuance de quantité, mais \eng{any} est la réponse attendue avec la négation simple. « Il ne reste pas de riz dans le sac. »

\item \eng{\textbf{Each} trader in the market pays a daily tax.} — Le nom \eng{trader} est au \textbf{singulier} et le verbe \eng{pays} est au singulier : seul \cle{each} convient. « Chaque commerçant du marché paie une taxe quotidienne. »

\item \eng{How \textbf{many} students passed the English test?} — \eng{Students} est un pluriel comptable : \cle{how many}. \eng{How much} est réservé aux incomptables (\eng{How much money?}). « Combien d'élèves ont réussi le test d'anglais ? »

\item \eng{We have very \textbf{little} time before the bus leaves.} — \eng{Time} est incomptable et le sens est négatif (« nous n'avons presque plus de temps ») : \cle{little} sans article. \eng{A little} signifierait « un peu de temps », ce qui contredit l'urgence.

\item \eng{She earns \textbf{little} money selling vegetables, but she is happy.} — \eng{Money} est incomptable ; le sens est « peu d'argent » (nuance négative soulignée par \eng{but}) : \cle{little}. « Elle gagne peu d'argent en vendant des légumes, mais elle est heureuse. »

\item \eng{There are \textbf{no} chairs in the classroom.} — Le verbe est déjà affirmatif (\eng{are}) et le sens est « il n'y a pas de chaises » : \cle{no} + nom, sans \eng{not}. On ne dit jamais \eng{There aren't no chairs} (double négation interdite en anglais standard).
\end{enumerate}

\textbf{Conclusion.} Le choix du quantifieur repose toujours sur trois questions : comptable ou incomptable ? phrase affirmative, négative ou interrogative ? sens positif ou négatif ?
\end{exoresolu}

\courssub{Méthode 2 : Éviter les erreurs fréquentes des francophones}

\begin{methode}[Les pièges du calque du français]
Le français et l'anglais ne répartissent pas les quantifieurs de la même façon. Démarche de vérification avant de rédiger :

\begin{enumerate}
\item \cle{Repérer le calque} : est-ce que je traduis mot à mot « beaucoup de », « peu de », « chaque », « tout » ?
\item \cle{Appliquer la règle anglaise}, pas la française :
\begin{itemize}
\item « beaucoup de » = \eng{a lot of} (registre courant) ; \eng{much/many} surtout en négatif et interrogatif ;
\item « chaque » = \eng{each} ou \eng{every} + \textbf{singulier} (jamais \eng{every students}) ;
\item « tous les deux » = \eng{both} (jamais \eng{the both}) ;
\item « aucun » = \eng{no} + nom ou \eng{none} seul (jamais de double négation) ;
\item « la plupart des » = \eng{most of the} + nom, ou \eng{most} + nom sans article (\eng{most students}) ; jamais \eng{the most students} au sens de « la plupart » (\eng{the most} = superlatif, « le plus »).
\end{itemize}
\item \cle{Vérifier l'accord du verbe} : \eng{each of them is}, \eng{none of the money was}, \eng{both of them are}.
\item \cle{Vérifier la place de \eng{of}} : \eng{some of my friends}, \eng{a few of the books}, mais \eng{some friends}, \eng{a few books} (sans déterminant, pas de \eng{of}).
\end{enumerate}
\end{methode}

\begin{exoresolu}[Exercice type : correction de phrases fautives]
\textbf{Énoncé.} Each sentence contains one mistake made by French-speaking learners. Correct it.

\begin{enumerate}
\item Every students in Form Five must wear a uniform.
\item I don't have no money for the bendskin.
\item The both brothers work at the port of Douala.
\item She gave me many advices about the job interview.
\item Most of students prefer English to German.
\item How much people attended the trade fair in Yaoundé?
\end{enumerate}

\tcblower
\textbf{Solution détaillée.}

\begin{enumerate}
\item \eng{Every students} $\rightarrow$ \eng{\textbf{Every student} in Form Five must wear a uniform.} — Après \cle{every} (comme après \cle{each}), le nom est toujours au \textbf{singulier}. Calque du français « tous les élèves ». « Tous les élèves de Première doivent porter l'uniforme. »

\item \eng{I don't have no money} $\rightarrow$ \eng{I \textbf{don't have any} money} ou \eng{I have \textbf{no} money for the bendskin.} — L'anglais standard interdit la \textbf{double négation}. Une seule négation par phrase : soit sur le verbe (\eng{don't \dots{} any}), soit sur le déterminant (\eng{no}). « Je n'ai pas d'argent pour le bendskin. »

\item \eng{The both brothers} $\rightarrow$ \eng{\textbf{Both} brothers work at the port of Douala} ou \eng{\textbf{Both of the} brothers work at the port of Douala.} — \cle{Both} ne prend jamais l'article \eng{the} devant lui. Avec \eng{of}, l'article devient possible : \eng{both of the brothers}. « Les deux frères travaillent au port de Douala. »

\item \eng{many advices} $\rightarrow$ \eng{She gave me \textbf{a lot of advice} about the job interview.} — \eng{Advice} est \textbf{incomptable} en anglais (faux ami du français « des conseils ») : jamais de \eng{-s}, jamais \eng{many}. Pour compter : \eng{a piece of advice}, \eng{two pieces of advice}. « Elle m'a donné beaucoup de conseils pour l'entretien d'embauche. »

\item \eng{Most of students} $\rightarrow$ \eng{\textbf{Most students} prefer English to German} ou \eng{\textbf{Most of the} students prefer English to German.} — \eng{Most of} exige un déterminant après lui (\eng{the, my, these}). Sans déterminant, on emploie \eng{most} seul : \eng{most students}. « La plupart des élèves préfèrent l'anglais à l'allemand. »

\item \eng{How much people} $\rightarrow$ \eng{How \textbf{many} people attended the trade fair in Yaoundé?} — \eng{People} est un pluriel comptable : \cle{how many}. \eng{How much} est réservé aux incomptables (\eng{how much money, how much time}). « Combien de personnes ont visité la foire commerciale de Yaoundé ? »
\end{enumerate}

\textbf{Conclusion.} Avant de valider une phrase, vérifier : singularité après \eng{each/every}, unicité de la négation, comptabilité du nom, présence ou absence de \eng{of}.
\end{exoresolu}

\courssub{Méthode 3 : Employer les quantifieurs dans un texte rédigé}

\begin{methode}[Utiliser la structure dans des phrases et de courts textes]
Pour rédiger un paragraphe (description d'un marché, d'un métier, d'une journée de travail) en mobilisant les quantifieurs :

\begin{enumerate}
\item \cle{Planifier le contenu} : quelles quantités vais-je exprimer ? (argent, temps, clients, produits, travailleurs).
\item \cle{Varier les quantifieurs} : ne pas répéter \eng{a lot of} partout ; alterner \eng{many, much, a few, a little, most, some, no, all, both, each}.
\item \cle{Respecter la grammaire} dans chaque phrase : comptabilité, accord du verbe, place de \eng{of}.
\item \cle{Relire} en traquant les cinq pièges : \eng{every} + singulier, pas de double négation, \eng{advice/information/money} incomptables, \eng{both} sans \eng{the}, \eng{most of the} avec article.
\end{enumerate}

\textbf{Structures utiles à réutiliser :}
\begin{itemize}
\item \eng{Most of the traders sell food, but a few of them sell clothes.} « La plupart des commerçants vendent de la nourriture, mais quelques-uns vendent des vêtements. »
\item \eng{Both my parents work, and each of them earns a salary.} « Mes deux parents travaillent, et chacun gagne un salaire. »
\item \eng{There is little unemployment in this sector, so many young people apply.} « Il y a peu de chômage dans ce secteur, donc beaucoup de jeunes postulent. »
\item \eng{None of the customers complained about the prices.} « Aucun des clients ne s'est plaint des prix. »
\end{itemize}
\end{methode}

\begin{exoresolu}[Exercice type : composition guidée]
\textbf{Énoncé.} Write a paragraph of about 80 words describing economic activity in a market you know (for example Mokolo Market or Marché Central). Use at least \textbf{six} different determiners or quantifiers from the lesson and underline them.

\tcblower
\textbf{Solution détaillée.}

\textbf{Étape 1 — Plan.} Lieu : Mokolo Market, Yaoundé. Idées : beaucoup de commerçants (\eng{many}), chaque étal (\eng{each}), la plupart des clients (\eng{most}), un peu de poisson frais (\eng{a little}), peu de temps libre (\eng{little}), aucun jour de repos (\eng{no}), tous les deux (\eng{both}).

\textbf{Étape 2 — Rédaction.}

\eng{Mokolo Market is one of the busiest places in Yaoundé. \textbf{Many} traders arrive before six o'clock, and \textbf{each} of them prepares a stall. \textbf{Most of the} customers come in the morning, when there is still \textbf{a little} fresh fish and meat. During the day, the traders have \textbf{little} free time because \textbf{all} the stalls are full of buyers. \textbf{Both} my aunt and her neighbour sell tomatoes there; they have \textbf{no} day off, not even on Sundays.}

« Le marché Mokolo est l'un des endroits les plus animés de Yaoundé. De nombreux commerçants arrivent avant six heures, et chacun d'eux prépare son étal. La plupart des clients viennent le matin, quand il reste encore un peu de poisson et de viande frais. Pendant la journée, les commerçants ont peu de temps libre car tous les étals sont pleins d'acheteurs. Ma tante et sa voisine y vendent toutes les deux des tomates ; elles n'ont aucun jour de repos, pas même le dimanche. »

\textbf{Étape 3 — Vérification grammaticale.}
\begin{itemize}
\item \eng{many traders} : pluriel comptable, correct ;
\item \eng{each of them prepares} : verbe au singulier après \eng{each of}, correct ;
\item \eng{most of the customers} : \eng{of} + article \eng{the}, correct ;
\item \eng{a little fresh fish} : incomptable, sens positif « un peu », correct ;
\item \eng{little free time} : incomptable, sens négatif « peu de », correct ;
\item \eng{all the stalls are} : pluriel, verbe au pluriel, correct ;
\item \eng{both my aunt and her neighbour} : pas de \eng{the} devant \eng{both}, correct ;
\item \eng{no day off} : négation portée par \eng{no}, verbe affirmatif, correct.
\end{itemize}

\textbf{Conclusion.} Le paragraphe compte environ 80 mots et emploie huit quantifieurs variés et correctement construits.
\end{exoresolu}

\courssub{Méthode 4 : Transformer des phrases avec les quantifieurs}

\begin{methode}[Réussir les exercices de transformation]
Les transformations classiques portent sur les paires de sens opposé ou équivalent. Démarche :

\begin{enumerate}
\item \cle{Identifier la paire logique} : \eng{some} $\leftrightarrow$ \eng{any/no/none} ; \eng{a few} $\leftrightarrow$ \eng{few} ; \eng{all} $\leftrightarrow$ \eng{none/no} ; \eng{each} $\leftrightarrow$ \eng{every}.
\item \cle{Transformer le verbe si nécessaire} : passer de la forme affirmative à la forme négative (ou l'inverse) quand on change de quantifieur : \eng{There is no milk} = \eng{There isn't any milk}.
\item \cle{Conserver le sens} : relire la phrase transformée et comparer avec l'originale.
\item \cle{Contrôler l'accord} : singulier/pluriel du nom et du verbe après transformation.
\end{enumerate}

\textbf{Équivalences à mémoriser :}
\begin{center}
\begin{tabularx}{\linewidth}{|X|X|}
\hline
\rowcolor{popblueL} Phrase affirmative & Équivalent négatif ou voisin \\
\hline
\eng{She has some money.} & \eng{She doesn't have any money.} / \eng{She has no money.} \\
\hline
\eng{All the students passed.} & \eng{None of the students failed.} \\
\hline
\eng{There are a few shops open.} & \eng{There aren't many shops open.} \\
\hline
\eng{Each worker received a bonus.} & \eng{Every worker received a bonus.} \\
\hline
\end{tabularx}
\end{center}
\end{methode}

\begin{exoresolu}[Exercice type : transformation de phrases]
\textbf{Énoncé.} Rewrite each sentence using the word in brackets without changing the meaning.

\begin{enumerate}
\item There isn't any sugar left. (no)
\item All the candidates answered the questions. (none \dots{} failed)
\item She has a little money in her mobile account. (not much)
\item Every employee received a uniform. (each)
\item Both of the banks are closed on Sundays. (neither \dots{} open)
\end{enumerate}

\tcblower
\textbf{Solution détaillée.}

\begin{enumerate}
\item \eng{There isn't any sugar left.} $\rightarrow$ \eng{There is \textbf{no} sugar left.} — On supprime la négation du verbe (\eng{isn't any}) et on la transfère sur le déterminant \eng{no}. Une seule négation dans la phrase. « Il ne reste pas de sucre. »

\item \eng{All the candidates answered the questions.} $\rightarrow$ \eng{\textbf{None of the candidates failed} to answer the questions.} — Totalité positive = absence d'échec. On construit \eng{fail to do something} (« ne pas réussir à faire »). « Tous les candidats ont répondu = aucun n'a échoué à répondre. »

\item \eng{She has a little money in her mobile account.} $\rightarrow$ \eng{She \textbf{doesn't have much} money in her mobile account.} — \eng{A little} (petite quantité positive) équivaut à \eng{not much}. \eng{Money} est incomptable : \eng{much} et non \eng{many}. Auxiliaire \eng{does} + base verbale \eng{have}. « Elle a un peu d'argent = elle n'a pas beaucoup d'argent. »

\item \eng{Every employee received a uniform.} $\rightarrow$ \eng{\textbf{Each} employee received a uniform.} — \eng{Each} et \eng{every} sont interchangeables ici ; tous deux exigent le nom au singulier et le verbe au singulier. Nuance : \eng{each} individualise davantage, \eng{every} généralise.

\item \eng{Both of the banks are closed on Sundays.} $\rightarrow$ \eng{\textbf{Neither of the banks is} open on Sundays.} — « Toutes les deux sont fermées » = « aucune des deux n'est ouverte ». Attention : \eng{neither of the banks} prend le verbe au \textbf{singulier} (\eng{is}) en anglais soigné.
\end{enumerate}

\textbf{Conclusion.} Dans une transformation, deux choses changent toujours ensemble : le quantifieur et la polarité (affirmative/négative) de la phrase. Vérifier enfin l'accord du verbe.
\end{exoresolu}

\begin{attention}[Les 5 erreurs qui coûtent des points au Bac]
\begin{enumerate}
\item \textbf{Le pluriel après \eng{each} et \eng{every}.} Écrire \eng{every students} ou \eng{each children} est une faute éliminatoire en grammaire : toujours le \textbf{singulier} — \eng{every student}, \eng{each child}, avec un verbe au singulier (\eng{every student \textbf{works}}).

\item \textbf{La double négation.} \eng{I don't have no money} est un calque du français « je n'ai pas d'argent ». Une seule négation : \eng{I don't have \textbf{any} money} ou \eng{I have \textbf{no} money}.

\item \textbf{Les faux amis incomptables.} \eng{Advice}, \eng{information}, \eng{news}, \eng{money}, \eng{work} sont \textbf{incomptables} : jamais de \eng{-s}, jamais \eng{many}. On écrit \eng{a piece of advice}, \eng{much information}, \eng{a lot of money}.

\item \textbf{\eng{The both} et \eng{most of students}.} \eng{Both} ne prend jamais \eng{the} devant : \eng{both brothers} ou \eng{both of the brothers}. \eng{Most of} exige un déterminant : \eng{most of \textbf{the} students} ; sinon \eng{most students} sans \eng{of}.

\item \textbf{Confondre \eng{how much} et \eng{how many}.} Devant un pluriel comptable, c'est \eng{how many} : \eng{How \textbf{many} people?}, \eng{How \textbf{many} francs?}. \eng{How much} est réservé aux incomptables (\eng{How much time?}) et au prix : \eng{How much is this shirt?} (« Combien coûte cette chemise ? »).
\end{enumerate}
\end{attention}

\newpage

\coursec{Exercices}

\courssub{Je vérifie mes connaissances}

\begin{exercice}[QCM : le bon quantifieur]
Entoure la bonne réponse.
\begin{enumerate}
\item There isn't \_\_\_\_\_ milk in the fridge. (some / any / no / many)
\item \_\_\_\_\_ students in our class passed the test. It was a success! (Most / Much / Any / None)
\item How \_\_\_\_\_ money do you need for the taxi? (many / much / few / any)
\item I have very \_\_\_\_\_ friends in Bamenda, only two or three. (little / much / few / any)
\item \_\_\_\_\_ of my parents works in Douala: my father is a trader and my mother is a nurse. (Each / Every / Both / All)
\item There is \_\_\_\_\_ water left in the bottle; we must buy some more. (a little / little / few / many)
\end{enumerate}
\end{exercice}

\begin{exercice}[Vrai ou faux ? Justifie]
Dis si les phrases suivantes sont correctes (\emph{right}) ou incorrectes (\emph{wrong}). Corrige les phrases fausses.
\begin{enumerate}
\item I have much friends at school.
\item She doesn't have any money today.
\item Every students are present in the classroom.
\item None of the shops was open on Sunday.
\item I don't have no time to help you.
\item Both of my brothers live in Yaoundé.
\end{enumerate}
\end{exercice}

\begin{exercice}[Vocabulaire : dénombrable ou indénombrable ?]
Classe les noms suivants en deux colonnes : \emph{countable nouns} / \emph{uncountable nouns}. Puis écris une phrase avec \eng{many} pour un nom dénombrable et une phrase avec \eng{much} pour un nom indénombrable.

\eng{money -- bread -- customer -- information -- shop -- rice -- job -- advice -- market -- water -- product -- news}
\end{exercice}

\begin{exercice}[Complète avec how much ou how many]
Complète les questions avec \eng{how much} ou \eng{how many}.
\begin{enumerate}
\item \_\_\_\_\_ does this fabric cost? -- Two thousand francs a metre.
\item \_\_\_\_\_ subjects do you study in Form Five?
\item \_\_\_\_\_ time do we have before the exam?
\item \_\_\_\_\_ of your friends work after school?
\item \_\_\_\_\_ sugar do you want in your tea?
\item \_\_\_\_\_ people came to the market this morning?
\end{enumerate}
\end{exercice}

\courssub{Je m'entraîne}

\begin{exercice}[Traduction français $\rightarrow$ anglais]
Traduis les phrases suivantes en anglais.
\begin{enumerate}
\item Combien coûte ce tissu ?
\item Il reste peu de tissu, mais quelques pièces sont en promotion.
\item La plupart des clients paient en liquide.
\item Je n'ai pas d'argent sur moi aujourd'hui.
\item Chaque commerçant du marché paie une taxe.
\item Aucun de mes amis n'a de compte en banque.
\end{enumerate}
\end{exercice}

\begin{exercice}[Transformation : négatif et interrogatif]
Mets chaque phrase à la forme négative, puis à la forme interrogative, en adaptant les quantifieurs.

Exemple : \eng{She has some money.} $\rightarrow$ \eng{She doesn't have any money.} $\rightarrow$ \eng{Does she have any money?}
\begin{enumerate}
\item There are some customers in the shop.
\item He bought some vegetables at the market.
\item They have some information about the new bank.
\item We saw some traders near the post office.
\end{enumerate}
\end{exercice}

\begin{exercice}[Texte à trous : au marché de Mokolo]
Complète le texte avec : \eng{some -- any -- much -- many -- few -- a few -- little -- a little -- all -- most -- each -- every} (chaque mot une seule fois).

\begin{textebox}[At Mokolo market]
\_\_\_\_\_ morning, hundreds of traders arrive at Mokolo market in Yaoundé. \_\_\_\_\_ of them sell food, but \_\_\_\_\_ sell clothes or shoes. There are so \_\_\_\_\_ customers on Saturdays that there is very \_\_\_\_\_ space to walk. ``I don't have \_\_\_\_\_ time to rest,'' says Mama Béa, a tomato seller. ``\_\_\_\_\_ customer wants a good price, and I sell \_\_\_\_\_ my tomatoes before noon.'' She earns \_\_\_\_\_ money every day, but not \_\_\_\_\_ : just enough to feed her family. ``\_\_\_\_\_ traders here are women,'' she adds. ``We work hard, and we save \_\_\_\_\_ money every week in a tontine.''
\end{textebox}
\end{exercice}

\begin{exercice}[Appariement : questions et réponses]
Associe chaque question (1--6) à la réponse qui convient (a--f).
\begin{enumerate}
\item How much is a bendskin ride to the post office?
\item How many workers are there in your father's shop?
\item Do you have any change?
\item How much homework do you have tonight?
\item Are there any banks near your school?
\item How many of your classmates have a part-time job?
\end{enumerate}
\begin{itemize}
\item[a.] \eng{None of them. We are all full-time students.}
\item[b.] \eng{A lot! Two essays and some maths exercises.}
\item[c.] \eng{Sorry, I have no coins at all.}
\item[d.] \eng{Only a few -- three, I think.}
\item[e.] \eng{Yes, there are two on the main road.}
\item[f.] \eng{About two hundred francs.}
\end{itemize}
\end{exercice}

\courssub{Je me prépare au Bac}

\begin{exercice}[Type Bac : compréhension et langue]
Lis le texte suivant, puis réponds aux questions.

\begin{textebox}[Small traders, big hopes]
In most Cameroonian towns, small businesses are everywhere: a tailor under a veranda, a woman selling puff-puff near a school, a young man repairing phones in a wooden kiosk. Many of these traders have little money to start with, but all of them have big dreams.

``When I opened my shop, I had very few customers,'' explains Jean, a barber in Bafoussam. ``Some days, I earned almost no money at all. But I didn't give up. Now I have many customers, and each of them pays five hundred francs for a haircut.''

Not every story is a success, however. A few traders borrow too much money and cannot pay it back. ``Most of us save a little money every week,'' says Rose, who sells second-hand clothes. ``It is the only way to survive when business is slow.'' None of these traders went to university, but both of the ones we interviewed said the same thing: hard work and patience matter more than diplomas.
\end{textebox}

\textbf{A. Compréhension.} Réponds en anglais par des phrases complètes.
\begin{enumerate}
\item Where do small traders usually work in Cameroonian towns? Give two examples from the text.
\item What problem did Jean have when he opened his shop?
\item According to Rose, how do traders survive when business is slow?
\item What quality matters more than diplomas, according to the two traders?
\end{enumerate}

\textbf{B. Langue.}
\begin{enumerate}
\item Trouve dans le texte trois quantifieurs différents et recopie la phrase qui les contient.
\item Mets à la forme négative : \eng{``Some days, I earned some money.''}
\item Pose la question correspondant à la réponse soulignée : \eng{``Each customer pays \underline{five hundred francs}.''}
\item Corrige l'erreur : \eng{``The most of traders save money every week.''}
\end{enumerate}
\end{exercice}

\begin{exercice}[Type Bac : traduction]
Traduis le passage suivant en anglais (5 phrases).

« Dans la plupart des villages, peu de gens ont un compte en banque. Chaque famille cultive un peu de maïs et quelques légumes. Combien d'argent gagnent-elles ? Très peu, mais aucune famille ne manque de nourriture. Beaucoup de jeunes partent en ville parce qu'il n'y a pas de travail au village. »
\end{exercice}

\begin{exercice}[Type Bac : rédaction guidée]
Rédige un paragraphe d'environ \textbf{10 lignes} (80--100 mots) intitulé \eng{``Small businesses in my town''}.

Consignes obligatoires :
\begin{itemize}
\item utilise les six quantifieurs suivants (souligne-les dans ta copie) : \eng{most, both, a few, no, every, many} ;
\item utilise au moins deux noms indénombrables avec \eng{much} ou \eng{little} ;
\item termine par une phrase donnant ton opinion personnelle.
\end{itemize}

Exemple de début : \eng{``In my town, most small businesses are shops and food stalls...''}
\end{exercice}

\newpage

\coursec{Corrigés des exercices}

\begin{corrige}[Exercice 1 — QCM : le bon quantifieur]
\begin{enumerate}
\item \eng{There isn’t \textbf{any} milk in the fridge.} — \eng{any} dans une phrase négative devant un nom indénombrable (\eng{milk}). « Il n’y a pas de lait dans le frigo. »
\item \eng{\textbf{Most} students in our class passed the test.} — \eng{most} + nom pluriel sans \eng{of} (généralité). « La plupart des élèves de notre classe ont réussi le test. »
\item \eng{How \textbf{much} money do you need for the taxi?} — \eng{money} est indénombrable $\rightarrow$ \eng{how much}. « De combien d’argent as-tu besoin pour le taxi ? »
\item \eng{I have very \textbf{few} friends in Bamenda, only two or three.} — \eng{friends} est dénombrable pluriel ; \eng{few} exprime la petitesse du nombre (valeur négative = « presque pas »). « J’ai très peu d’amis à Bamenda. »
\item \eng{\textbf{Both} of my parents \textbf{work} in Douala.} — \eng{both of} + nom pluriel + verbe au pluriel. « Mes deux parents travaillent à Douala. »
\item \eng{There is \textbf{little} water left in the bottle.} — \eng{little} (valeur négative : « presque plus ») justifié par la suite : \eng{we must buy some more}. « Il reste peu d’eau dans la bouteille. »
\end{enumerate}
\end{corrige}

\begin{corrige}[Exercice 2 — Vrai ou faux ? Justifie]
\begin{enumerate}
\item \emph{Wrong.} \eng{I have \textbf{many} friends at school.} — \eng{friends} est dénombrable ; \eng{much} ne s’emploie qu’avec les indénombrables.
\item \emph{Right.} \eng{any} + indénombrable (\eng{money}) dans une phrase négative : correct.
\item \emph{Wrong.} \eng{\textbf{Every student is} present in the classroom.} — \eng{every} est toujours suivi d’un nom au \textbf{singulier} et d’un verbe au singulier.
\item \emph{Right.} \eng{none of} + nom pluriel peut prendre un verbe au singulier (style soigné) ou au pluriel (\eng{were} accepté à l’oral).
\item \emph{Wrong.} \eng{I don’t have \textbf{any} time to help you.} (ou \eng{I have \textbf{no} time to help you.}) — jamais de double négation en anglais : \eng{not} + \eng{no} est impossible.
\item \emph{Right.} \eng{both of} + nom pluriel + verbe au pluriel : correct. « Mes deux frères habitent à Yaoundé. »
\end{enumerate}
\end{corrige}

\begin{corrige}[Exercice 3 — Vocabulaire : dénombrable ou indénombrable ?]
\begin{center}
\begin{tabularx}{\linewidth}{|X|X|}
\hline
\rowcolor{popblueL} \textbf{Countable nouns} & \textbf{Uncountable nouns} \\
\hline
customer, shop, job, market, product & money, bread, information, rice, advice, water, news \\
\hline
\end{tabularx}
\end{center}
Phrases possibles :
\begin{itemize}
\item \eng{There are \textbf{many} customers at the market on Saturdays.} « Il y a beaucoup de clients au marché le samedi. »
\item \eng{She doesn’t earn \textbf{much} money in her small shop.} « Elle ne gagne pas beaucoup d’argent dans sa petite boutique. »
\end{itemize}
\textbf{Pièges à retenir :} \eng{information}, \eng{advice}, \eng{news} et \eng{money} sont indénombrables en anglais alors que leurs équivalents français prennent souvent un pluriel (« des informations », « des conseils »). On ne dit jamais \eng{advices} ni \eng{informations}.
\end{corrige}

\begin{corrige}[Exercice 4 — Complète avec \eng{how much} ou \eng{how many}]
\begin{enumerate}
\item \eng{\textbf{How much} does this fabric cost?} — on demande un prix (indénombrable).
\item \eng{\textbf{How many} subjects do you study in Form Five?} — \eng{subjects} : dénombrable pluriel.
\item \eng{\textbf{How much} time do we have before the exam?} — \eng{time} (la durée) est indénombrable.
\item \eng{\textbf{How many} of your friends work after school?} — \eng{how many of} + pronom/nom défini pluriel.
\item \eng{\textbf{How much} sugar do you want in your tea?} — \eng{sugar} : indénombrable.
\item \eng{\textbf{How many} people came to the market this morning?} — \eng{people} : pluriel dénombrable.
\end{enumerate}
\textbf{Règle :} \eng{how much} + indénombrable ; \eng{how many} + dénombrable pluriel. Ne jamais dire \eng{how many money}.
\end{corrige}

\begin{corrige}[Exercice 5 — Traduction français $\rightarrow$ anglais]
\begin{enumerate}
\item \eng{How much does this fabric cost?} (ou \eng{How much is this fabric?})
\item \eng{There is \textbf{little} fabric left, but \textbf{a few} pieces are on sale.} — \eng{little} + indénombrable ; \eng{a few} + pluriel.
\item \eng{\textbf{Most} customers pay \textbf{in cash}.} (ou \eng{Most of the customers pay cash.}) — \eng{in cash}, pas \eng{in liquid} (faux ami) !
\item \eng{I don’t have \textbf{any} money on me today.} (ou \eng{I have \textbf{no} money on me today.})
\item \eng{\textbf{Each} trader in the market pays a tax.} — \eng{each} + singulier ; \eng{every trader} serait aussi correct.
\item \eng{\textbf{None} of my friends \textbf{has} a bank account.} (ou \eng{have}, accepté) — \eng{none of} + nom pluriel ; pas de double négation.
\end{enumerate}
\end{corrige}

\begin{corrige}[Exercice 6 — Transformation : négatif et interrogatif]
\begin{enumerate}
\item \eng{There are \textbf{some} customers in the shop.} $\rightarrow$ \eng{There \textbf{aren’t any} customers in the shop.} $\rightarrow$ \eng{\textbf{Are there any} customers in the shop?}
\item \eng{He bought \textbf{some} vegetables at the market.} $\rightarrow$ \eng{He \textbf{didn’t buy any} vegetables at the market.} $\rightarrow$ \eng{\textbf{Did he buy any} vegetables at the market?}
\item \eng{They have \textbf{some} information about the new bank.} $\rightarrow$ \eng{They \textbf{don’t have any} information about the new bank.} $\rightarrow$ \eng{\textbf{Do they have any} information about the new bank?}
\item \eng{We saw \textbf{some} traders near the post office.} $\rightarrow$ \eng{We \textbf{didn’t see any} traders near the post office.} $\rightarrow$ \eng{\textbf{Did you see any} traders near the post office?} (sujet changé en \eng{you} pour la question naturelle).
\end{enumerate}
\textbf{Règle appliquée :} \eng{some} devient \eng{any} à la forme négative et à la forme interrogative ; attention au prétérit : l’auxiliaire \eng{did} reprend la marque du passé (\eng{bought} $\rightarrow$ \eng{didn’t buy}).
\end{corrige}

\begin{corrige}[Exercice 7 — Texte à trous : au marché de Mokolo]
Texte complété :

\eng{\textbf{Every} morning, hundreds of traders arrive at Mokolo market in Yaoundé. \textbf{Most} of them sell food, but \textbf{some} sell clothes or shoes. There are so \textbf{many} customers on Saturdays that there is very \textbf{little} space to walk. “I don’t have \textbf{any} time to rest,” says Mama Béa, a tomato seller. “\textbf{Each} customer wants a good price, and I sell \textbf{all} my tomatoes before noon.” She earns \textbf{a little} money every day, but not \textbf{much}: just enough to feed her family. “\textbf{Most} traders here are women… we save \textbf{some} money every week.”}

\textbf{Attention :} le texte contient douze blancs mais deux d’entre eux (\eng{Most of them} et \eng{Most traders}) utilisent le même mot ; la consigne « chaque mot une seule fois » s’applique donc avec cette réserve. Répartition justifiée :
\begin{itemize}
\item \eng{Every morning} : \eng{every} + singulier pour une habitude.
\item \eng{Most of them} / \eng{some sell} : opposition « la plupart / certains ».
\item \eng{so many customers} : dénombrable pluriel ; \eng{very little space} : indénombrable, valeur négative.
\item \eng{any time} : phrase négative ; \eng{Each customer} : singulier, insiste sur l’individu.
\item \eng{all my tomatoes} : totalité ; \eng{a little money} : « un peu » (positif) ; \eng{not much} : « pas beaucoup ».
\item Dernier blanc : \eng{some money} (indénombrable, valeur positive) — \eng{a few} impossible devant \eng{money}.
\end{itemize}
\end{corrige}

\begin{corrige}[Exercice 8 — Appariement : questions et réponses]
\begin{center}
\begin{tabularx}{\linewidth}{|l|X|X|}
\hline
\rowcolor{popblueL} \textbf{N°} & \textbf{Réponse} & \textbf{Justification} \\
\hline
1 & f & prix d’une course : “About two hundred francs.” \\
\hline
2 & d & nombre de travailleurs : “Only a few -- three, I think.” \\
\hline
3 & c & monnaie : “Sorry, I have no coins at all.” \\
\hline
4 & b & quantité de devoirs (indénombrable) : “A lot!” \\
\hline
5 & e & existence de banques : “Yes, there are two...” \\
\hline
6 & a & \eng{how many of} $\rightarrow$ “None of them.” \\
\hline
\end{tabularx}
\end{center}
\end{corrige}

\begin{corrige}[Exercice 9 — Type Bac : compréhension et langue]
\textbf{A. Compréhension} (réponses modèles)
\begin{enumerate}
\item \eng{Small traders usually work in simple places in the street: for example, a tailor works under a veranda and a woman sells puff-puff near a school.}
\item \eng{When he opened his shop, he had very few customers, and some days he earned almost no money at all.}
\item \eng{According to Rose, they survive by saving a little money every week.}
\item \eng{According to the two traders, hard work and patience matter more than diplomas.}
\end{enumerate}

\textbf{B. Langue}
\begin{enumerate}
\item Trois quantifieurs possibles (au choix) : \eng{\textbf{Many} of these traders have little money...} ; \eng{\textbf{All} of them have big dreams} ; \eng{\textbf{Each} of them pays five hundred francs} ; \eng{\textbf{Most} of us save a little money every week} ; \eng{\textbf{None} of these traders went to university}.
\item \eng{“Some days, I \textbf{didn’t earn any} money.”} (ou \eng{I earned no money.})
\item \eng{\textbf{How much} does each customer pay?} — on interroge sur un prix $\rightarrow$ \eng{how much}.
\item \eng{“\textbf{Most} traders save money every week.”} (ou \eng{\textbf{Most of the} traders...}) — jamais \eng{the most of}.
\end{enumerate}

\textbf{Barème indicatif (10 pts) :} A : 1 pt par réponse correcte et bien formulée (4 pts). B : 1,5 pt par item (6 pts). Points de vigilance : phrase complète exigée en compréhension ; orthographe de \eng{customers}, \eng{money}, \eng{business} ; pas de double négation.
\end{corrige}

\begin{corrige}[Exercice 10 — Type Bac : traduction]
Proposition de traduction :

\eng{“In most villages, few people have a bank account. Each family grows a little maize and a few vegetables. How much money do they earn? Very little, but no family lacks food. Many young people leave for the town because there is no work in the village.”}

Points de vigilance :
\begin{itemize}
\item « peu de gens » $\rightarrow$ \eng{few people} (dénombrable) ; « un peu de maïs » $\rightarrow$ \eng{a little maize} (indénombrable) ; « quelques légumes » $\rightarrow$ \eng{a few vegetables}.
\item « Combien d’argent » $\rightarrow$ \eng{how much money} (jamais \eng{how many money}).
\item « aucune famille » $\rightarrow$ \eng{no family} + verbe affirmatif (pas de double négation).
\item « beaucoup de jeunes » $\rightarrow$ \eng{many young people} ; \eng{people} sans \eng{s}.
\end{itemize}
\textbf{Barème indicatif (5 pts) :} 1 pt par phrase ; 0,5 si le sens est correct mais le quantifieur est fautif.
\end{corrige}

\begin{corrige}[Exercice 11 — Type Bac : rédaction guidée]
\textbf{Proposition de rédaction modèle} (les quantifieurs imposés sont en gras) :

\eng{“In my town, \textbf{most} small businesses are shops and food stalls. Near the market, there are \textbf{many} traders who sell fruit, vegetables and second-hand clothes. My uncle and my aunt run a small restaurant: \textbf{both} of them work from morning to night. \textbf{Every} day, they serve \textbf{a few} regular customers who come for lunch. They do not earn \textbf{much} \textbf{money}, but they never have \textbf{little} \textbf{work} to do. There is \textbf{no} bank in our quarter, so people keep their savings in a tontine. In my opinion, small businesses are the heart of our town because they give work and hope to everybody.”}

(environ 95 mots)

\textbf{Barème indicatif (10 pts) :} respect de la longueur et du sujet (2 pts) ; les huit quantifieurs imposés correctement employés et mis en gras (3 pts) ; deux indénombrables avec \eng{much}/\eng{little} (2 pts) ; opinion personnelle finale (1 pt) ; correction grammaticale et orthographe (2 pts).

Points de vigilance : \eng{most} sans \eng{the} devant un nom général ; \eng{both of them} + verbe pluriel ; \eng{every} + singulier ; \eng{no} + verbe affirmatif ; éviter les tournures calquées du français (\eng{there is no of work} $\rightarrow$ \eng{there is no work}).
\end{corrige}

\newpage

\coursec{Fiche bilan}

\begin{popfigure}
\begin{tikzpicture}[mindmap, grow cyclic, every node/.style={concept, align=center, font=\small},
root concept/.append style={concept color=popblue, font=\small\bfseries, minimum size=3.2cm},
level 1 concept/.append style={sibling angle=60, level distance=3.6cm, font=\scriptsize, minimum size=2.4cm}]
\node[root concept]{Determiners and quantifiers}
  child[concept color=poporange]{ node[align=center]{some / any / no / none\\(affirmatif, interrogatif, négatif)} }
  child[concept color=popgreen]{ node[align=center]{much / many\\how much / how many} }
  child[concept color=poppurple]{ node[align=center]{few / a few\\little / a little} }
  child[concept color=poppink]{ node[align=center]{all (of) / both (of)\\most (of)} }
  child[concept color=popgold]{ node[align=center]{each / every\\(singulier)} }
  child[concept color=popblue!70]{ node[align=center]{countable vs\\uncountable} };
\end{tikzpicture}
\legende{Carte mentale de la leçon : les déterminants et quantifieurs anglais.}
\end{popfigure}

\begin{bilan}
\textbf{Lexique clé (English / Français)}

\begin{center}
\begin{tabularx}{\linewidth}{|X|X|X|X|}
\hline
\rowcolor{popblueL} English & Français & English & Français \\
\hline
some & du, des, quelques & none (of) & aucun(e) (de) \\
\hline
any & du, de, aucun (int./nég.) & few & peu de (nombrables) \\
\hline
no & aucun, pas de & a few & quelques (nombrables) \\
\hline
much & beaucoup de (indénombr.) & little & peu de (indénombr.) \\
\hline
many & beaucoup de (dénombr.) & a little & un peu de (indénombr.) \\
\hline
all (of) & tout, tous, toutes & each & chaque (individuellement) \\
\hline
both (of) & les deux & every & chaque, tout (généralité) \\
\hline
most (of) & la plupart (de) & how much / how many & combien de \\
\hline
\end{tabularx}
\end{center}

\textbf{Règles à connaître par cœur}

\begin{center}
\begin{tabularx}{\linewidth}{|l|X|X|}
\hline
\rowcolor{popblueL} Quantifieur & Nom dénombrable (books) & Nom indénombrable (money) \\
\hline
some / any / no / none & some books, any books & some money, any money \\
\hline
much / many & many books & much money \\
\hline
few / a few & few books, a few books & --- \\
\hline
little / a little & --- & little money, a little money \\
\hline
all / most & all books, most books & all money, most money \\
\hline
each / every & each book, every book (singulier) & --- \\
\hline
\end{tabularx}
\end{center}

\begin{itemize}
\item \cle{some} s'emploie dans les phrases affirmatives : \eng{There is some rice on the table.} ; \cle{any} dans les phrases interrogatives et négatives : \eng{Do you have any money?} / \eng{I don't have any money.}
\item Exception : \cle{some} dans les questions qui sont des offres ou des demandes : \eng{Would you like some tea?} / \eng{Can I have some water?}
\item \cle{no} + nom = phrase affirmative au sens négatif : \eng{I have no money.} (= \eng{I don't have any money.}) ; jamais deux négations.
\item \cle{none (of)} remplace le nom : \eng{How many students passed? None.} / \eng{None of my friends came.}
\item \cle{much} + indénombrable ; \cle{many} + dénombrable pluriel. Questions : \eng{How much money?} / \eng{How many books?}
\item \cle{few} / \cle{little} = « peu de » (sens négatif, insuffisant) ; \cle{a few} / \cle{a little} = « quelques / un peu de » (sens positif) : \eng{Few people came.} (presque personne) vs \eng{A few people came.} (quelques personnes).
\item \cle{each} et \cle{every} sont suivis d'un nom au singulier et d'un verbe au singulier : \eng{Every student has a book.} ; \cle{each} insiste sur l'individu, \cle{every} sur le groupe.
\item \cle{all (of)}, \cle{both (of)}, \cle{most (of)} : \emph{of} est obligatoire devant un déterminant ou un pronom : \eng{All of the students}, \eng{Both of them}, mais \eng{All students}, \eng{Most people}.
\item \cle{both} = exactement deux : \eng{Both my parents are teachers.} ; verbe au pluriel.
\end{itemize}

\textbf{Méthodes express}

\begin{enumerate}
\item Identifier le nom : dénombrable (singulier/pluriel) ou indénombrable (\eng{money, water, information, advice, news, furniture}).
\item Identifier le type de phrase : affirmative (\cle{some}), négative ou interrogative (\cle{any}), négation par \cle{no}/\cle{none}.
\item Choisir le quantifieur selon le sens : quantité grande (\cle{much/many}), petite positive (\cle{a few/a little}), petite négative (\cle{few/little}), totalité (\cle{all/both/each/every/most}).
\item Vérifier l'accord du verbe : singulier après \cle{each/every}, pluriel après \cle{both/few/many}.
\end{enumerate}
\end{bilan}

\begin{aretenir}[Checklist avant le Bac]
\begin{itemize}
\item[$\square$] Je distingue les noms dénombrables des indénombrables (\eng{information, money, bread, news} = indénombrables).
\item[$\square$] J'utilise \cle{some} à l'affirmatif et \cle{any} à l'interrogatif et au négatif.
\item[$\square$] Je sais que \cle{some} reste possible dans les offres et demandes : \eng{Would you like some coffee?}
\item[$\square$] J'écris \eng{I have no money} ou \eng{I don't have any money}, jamais de double négation.
\item[$\square$] J'emploie \cle{much} avec les indénombrables et \cle{many} avec les dénombrables.
\item[$\square$] Je pose les questions de quantité avec \eng{How much...?} ou \eng{How many...?} correctement.
\item[$\square$] Je fais la différence entre \cle{few} (peu, sens négatif) et \cle{a few} (quelques, sens positif), \cle{little} et \cle{a little}.
\item[$\square$] Je mets le nom et le verbe au singulier après \cle{each} et \cle{every}.
\item[$\square$] Je mets \emph{of} après \cle{all, both, most, none} devant \eng{the, my, these} ou un pronom : \eng{Most of the students}.
\item[$\square$] J'emploie \cle{both} seulement pour deux éléments, avec un verbe au pluriel.
\item[$\square$] Je n'ajoute jamais de \emph{s} aux indénombrables : \eng{some information}, pas \eng{informations}.
\item[$\square$] Je sais répondre par \cle{none} à une question en \eng{How many...?} : \eng{None.}
\end{itemize}
\end{aretenir}

\findecours

\end{document}
